%% Generated by Sphinx.
\def\sphinxdocclass{report}
\documentclass[a4paper,11pt,english,openany,oneside]{sphinxmanual}
\ifdefined\pdfpxdimen
   \let\sphinxpxdimen\pdfpxdimen\else\newdimen\sphinxpxdimen
\fi \sphinxpxdimen=.75bp\relax
\ifdefined\pdfimageresolution
    \pdfimageresolution= \numexpr \dimexpr1in\relax/\sphinxpxdimen\relax
\fi
%% let collapsible pdf bookmarks panel have high depth per default
\PassOptionsToPackage{bookmarksdepth=5}{hyperref}

\PassOptionsToPackage{booktabs}{sphinx}
\PassOptionsToPackage{colorrows}{sphinx}

\PassOptionsToPackage{warn}{textcomp}
\usepackage[utf8]{inputenc}
\ifdefined\DeclareUnicodeCharacter
% support both utf8 and utf8x syntaxes
  \ifdefined\DeclareUnicodeCharacterAsOptional
    \def\sphinxDUC#1{\DeclareUnicodeCharacter{"#1}}
  \else
    \let\sphinxDUC\DeclareUnicodeCharacter
  \fi
  \sphinxDUC{00A0}{\nobreakspace}
  \sphinxDUC{2500}{\sphinxunichar{2500}}
  \sphinxDUC{2502}{\sphinxunichar{2502}}
  \sphinxDUC{2514}{\sphinxunichar{2514}}
  \sphinxDUC{251C}{\sphinxunichar{251C}}
  \sphinxDUC{2572}{\textbackslash}
\fi
\usepackage{cmap}
\usepackage[T1]{fontenc}
\usepackage{amsmath,amssymb,amstext}
\usepackage{babel}



\usepackage{tgtermes}
\usepackage{tgheros}
\renewcommand{\ttdefault}{txtt}



\usepackage[Bjarne]{fncychap}
\usepackage{sphinx}
\sphinxsetup{hmargin={0.7in,0.7in}, vmargin={0.7in,0.7in}, verbatimwithframe=false}
\fvset{fontsize=auto}
\usepackage{geometry}


% Include hyperref last.
\usepackage{hyperref}
% Fix anchor placement for figures with captions.
\usepackage{hypcap}% it must be loaded after hyperref.
% Set up styles of URL: it should be placed after hyperref.
\urlstyle{same}


\usepackage{sphinxmessages}
\setcounter{tocdepth}{0}


\usepackage{babel}
\usepackage{graphicx}
\usepackage{hyperref}
\setcounter{tocdepth}{2}
\raggedbottom

% TWARDY RESET INDEKSU - zmusza LaTeXa do zignorowania komendy generującej spis
\renewcommand{\printindex}{}

% Zmniejszenie przestrzeni przed nagłówkami
\setlength{\parskip}{0pt plus 1pt}
\setlength{\parindent}{0pt}

% Zwiększenie wysokości headera
\setlength{\headheight}{14.49998pt}


\title{Sprawozdanie z Laboratorium: Bazy Danych}
\date{Jun 13, 2026}
\release{}
\author{Paweł Łoćwin}
\newcommand{\sphinxlogo}{\vbox{}}
\renewcommand{\releasename}{}
\makeindex
\begin{document}

\ifdefined\shorthandoff
  \ifnum\catcode`\=\string=\active\shorthandoff{=}\fi
  \ifnum\catcode`\"=\active\shorthandoff{"}\fi
\fi

\pagestyle{empty}
\sphinxmaketitle
\pagestyle{plain}
\sphinxtableofcontents
\pagestyle{normal}
\phantomsection\label{\detokenize{index::doc}}


\sphinxAtStartPar
\sphinxstylestrong{Autorzy projektu:} Paweł Łoćwin, Paweł Łosowski

\sphinxAtStartPar
\sphinxstylestrong{Wprowadzenie:}

\sphinxAtStartPar
Niniejsza dokumentacja stanowi kompletne sprawozdanie z laboratorium przedmiotu Bazy Danych. Projekt obejmuje całościową analizę procesów związanych z administracją i projektowaniem relacyjnych baz danych.

\sphinxAtStartPar
Dokumentacja została podzielona na pięć głównych sekcji:
\begin{enumerate}
\sphinxsetlistlabels{\arabic}{enumi}{enumii}{}{.}%
\item {} 
\sphinxAtStartPar
\sphinxstylestrong{Rozdział 1: Wstęp oraz linki} \textendash{} Wstęp do sprawozdania oraz linki.

\item {} 
\sphinxAtStartPar
\sphinxstylestrong{Rozdział 2: Badania literaturowe} \textendash{} Przegląd źródeł akademickich i naukowo\sphinxhyphen{}technicznych dotyczących baz danych, ich architektury, optymalizacji i bezpieczeństwa.

\item {} 
\sphinxAtStartPar
\sphinxstylestrong{Rozdział 3: Projektowanie bazy danych} \textendash{} Praktyczne modelowanie systemu zarządzania bibliotecznym katalogiem wypożyczeń. Sekcja obejmuje analizę wymagań, model konceptualny (notacja E\sphinxhyphen{}R) oraz model logiczny.

\item {} 
\sphinxAtStartPar
\sphinxstylestrong{Rozdział 4: Implementacja fizyczna i zasilanie bazy} \textendash{} Realizacja fizycznego schematu bazy danych poprzez skrypty DDL, mechanizmy importu danych oraz automatyzacja zasilania bazy.

\item {} 
\sphinxAtStartPar
\sphinxstylestrong{Rozdział 5: Zapytania do bazy danych} \textendash{} Zaawansowane funkcje SQL w Python, selekcje, agregacje, złączenia, operatory zbiorowe i podzapytania z pełną dokumentacją Sphinx.

\end{enumerate}

\sphinxstepscope


\chapter{Rozdział 1: Wstęp oraz linki}
\label{\detokenize{rozdzial_1/index:rozdzial-1-wstep-oraz-linki}}\label{\detokenize{rozdzial_1/index::doc}}\begin{quote}\begin{description}
\sphinxlineitem{Autorzy}\begin{enumerate}
\sphinxsetlistlabels{\arabic}{enumi}{enumii}{}{.}%
\item {} 
\sphinxAtStartPar
Paweł Łoćwin

\item {} 
\sphinxAtStartPar
Paweł Łosowski

\end{enumerate}

\end{description}\end{quote}


\section{Wstęp do sprawozdania}
\label{\detokenize{rozdzial_1/index:wstep-do-sprawozdania}}
\sphinxAtStartPar
Niniejsze sprawozdanie dokumentuje przebieg prac laboratoryjnych z przedmiotu Bazy Danych. Celem projektu jest praktyczne zastosowanie wiedzy z zakresu projektowania, wdrażania oraz obsługi systemów relacyjnych baz danych.

\sphinxAtStartPar
Praca obejmuje pełen cykl życia bazy: od analizy teoretycznej, poprzez stworzenie modelu pojęciowego i logicznego, aż po fizyczną implementację w środowiskach PostgreSQL i SQLite, a także realizację procesów wsadowego zasilania danymi.


\section{Wykaz repozytoriów (Linki)}
\label{\detokenize{rozdzial_1/index:wykaz-repozytoriow-linki}}
\sphinxAtStartPar
Projekt został podzielony na odpowiednie repozytoria w systemie kontroli wersji Git, zapewniając porządek między dokumentacją a fizycznymi plikami środowiska bazodanowego.


\subsection{1. Główne repozytorium sprawozdania (Dokumentacja Sphinx)}
\label{\detokenize{rozdzial_1/index:glowne-repozytorium-sprawozdania-dokumentacja-sphinx}}
\sphinxAtStartPar
Tutaj znajduje się struktura całej dokumentacji tekstowej, konfiguracja kompilatora Sphinx oraz historia zmian.
\begin{itemize}
\item {} 
\sphinxAtStartPar
\sphinxstylestrong{Link:} \sphinxhref{https://github.com/baguetedev/Bazy-Danych-1}{baguetedev/Bazy\sphinxhyphen{}Danych\sphinxhyphen{}1}%
\begin{footnote}[1]\sphinxAtStartFootnote
\sphinxnolinkurl{https://github.com/baguetedev/Bazy-Danych-1}
%
\end{footnote}

\end{itemize}


\subsection{2. Repozytorium z plikami projektowymi}
\label{\detokenize{rozdzial_1/index:repozytorium-z-plikami-projektowymi}}
\sphinxAtStartPar
W tym repozytorium umieszczono kody źródłowe, skrypty DDL dla silników PostgreSQL i SQLite, wygenerowany plik bazy oraz pliki CSV wykorzystywane do zasilania bazy danych.
\begin{itemize}
\item {} 
\sphinxAtStartPar
\sphinxstylestrong{Link WWW:} \sphinxhref{https://github.com/baguetedev/Bazy-Danych-pliki}{baguetedev/Bazy\sphinxhyphen{}Danych\sphinxhyphen{}pliki}%
\begin{footnote}[2]\sphinxAtStartFootnote
\sphinxnolinkurl{https://github.com/baguetedev/Bazy-Danych-pliki}
%
\end{footnote}

\item {} 
\sphinxAtStartPar
\sphinxstylestrong{Klonowanie SSH:} \sphinxcode{\sphinxupquote{git@github.com:baguetedev/Bazy\sphinxhyphen{}Danych\sphinxhyphen{}pliki.git}}

\end{itemize}


\subsection{3. Repozytoria reszty grupy}
\label{\detokenize{rozdzial_1/index:repozytoria-reszty-grupy}}
\sphinxAtStartPar
Poniżej znajdują się odnośniki do repozytoriów reszty grupy, które zostały zintegrowane z niniejszym sprawozdaniem w postaci submodułów w Rozdziale 2:
\begin{itemize}
\item {} 
\sphinxAtStartPar
\sphinxstylestrong{Grupa 1 (rozdzial\_1):} \sphinxcode{\sphinxupquote{https://github.com/karaskamil/Sprzet\sphinxhyphen{}dla\sphinxhyphen{}bazy\sphinxhyphen{}danych.git}}

\item {} 
\sphinxAtStartPar
\sphinxstylestrong{Grupa 2 (rozdzial\_2):} \sphinxcode{\sphinxupquote{https://github.com/Youarecheck/Bazy\_Danych\_Tematyczne\_Repo\_MK.git}}

\item {} 
\sphinxAtStartPar
\sphinxstylestrong{Grupa 3 (rozdzial\_3):} \sphinxcode{\sphinxupquote{https://github.com/pawlos1337/Bazy\sphinxhyphen{}danych\sphinxhyphen{}temat.git}}

\item {} 
\sphinxAtStartPar
\sphinxstylestrong{Grupa 4 (rozdzial\_4):} \sphinxcode{\sphinxupquote{https://github.com/OskarProgrammer/monitorowanie\_i\_diagnostyka.git}}

\item {} 
\sphinxAtStartPar
\sphinxstylestrong{Grupa 5 (rozdzial\_5):} \sphinxcode{\sphinxupquote{https://github.com/KMachoK/Tematyczne.git}}

\item {} 
\sphinxAtStartPar
\sphinxstylestrong{Grupa 6 (rozdzial\_6):} \sphinxcode{\sphinxupquote{https://github.com/domino0472/Partycjonowani\sphinxhyphen{}Danych}}

\item {} 
\sphinxAtStartPar
\sphinxstylestrong{Grupa 7 (rozdzial\_7):} \sphinxcode{\sphinxupquote{https://github.com/oski486/BazyDanych\sphinxhyphen{}Subject.git}}

\item {} 
\sphinxAtStartPar
\sphinxstylestrong{Grupa 8 (rozdzial\_8):} \sphinxcode{\sphinxupquote{https://github.com/Koko9077/Kopie\sphinxhyphen{}zapasowe\sphinxhyphen{}i\sphinxhyphen{}odzyskiwanie\sphinxhyphen{}danych.git}}

\end{itemize}

\sphinxstepscope


\chapter{Badania literaturowe}
\label{\detokenize{rozdzial_2/index:badania-literaturowe}}\label{\detokenize{rozdzial_2/index::doc}}
\sphinxstepscope


\section{Sprzęt dla bazy danych}
\label{\detokenize{rozdzial_2/rozdzial_21/index:sprzet-dla-bazy-danych}}\label{\detokenize{rozdzial_2/rozdzial_21/index::doc}}
\sphinxstepscope


\section{Konfiguracja bazy danych PostgreSQL}
\label{\detokenize{rozdzial_2/rozdzial_22/index:konfiguracja-bazy-danych-postgresql}}\label{\detokenize{rozdzial_2/rozdzial_22/index::doc}}

\subsection{Wstęp}
\label{\detokenize{rozdzial_2/rozdzial_22/index:wstep}}
\sphinxAtStartPar
Rozdział opisuje konfigurację połączenia z bazą danych \sphinxstylestrong{PostgreSQL} \textendash{} dojrzałym, open\sphinxhyphen{}source’owym systemem zarządzania relacyjnymi bazami danych (RDBMS), który wyróżnia się pełną zgodnością ze standardem SQL oraz silnym naciskiem na rozszerzalność i niezawodność.

\sphinxAtStartPar
\sphinxstylestrong{Wymagania:}
\begin{itemize}
\item {} 
\sphinxAtStartPar
dostęp do instancji PostgreSQL,

\item {} 
\sphinxAtStartPar
zmienne środowiskowe,

\item {} 
\sphinxAtStartPar
narzędzia projektowe (np. \sphinxcode{\sphinxupquote{psql}}, \sphinxcode{\sphinxupquote{pg\_isready}}).

\end{itemize}


\subsection{Podstawowe parametry}
\label{\detokenize{rozdzial_2/rozdzial_22/index:podstawowe-parametry}}
\sphinxAtStartPar
Każde połączenie z PostgreSQL wymaga:
\begin{itemize}
\item {} 
\sphinxAtStartPar
hosta i portu (domyślnie \sphinxcode{\sphinxupquote{localhost}} i \sphinxcode{\sphinxupquote{5432}}),

\item {} 
\sphinxAtStartPar
nazwy bazy danych,

\item {} 
\sphinxAtStartPar
użytkownika i hasła.

\end{itemize}

\sphinxAtStartPar
Dane wrażliwe (np. hasła) przechowuj w \sphinxstylestrong{zmiennych środowiskowych} \textendash{} to oddziela konfigurację od kodu.
Pamiętaj: zmienne środowiskowe nie są mechanizmem bezpieczeństwa; w produkcji uzupełnij je o dedykowane
zarządzanie sekretami (np. HashiCorp Vault, AWS Secrets Manager).


\subsection{Lokalizacja i struktura katalogów PostgreSQL}
\label{\detokenize{rozdzial_2/rozdzial_22/index:lokalizacja-i-struktura-katalogow-postgresql}}
\sphinxAtStartPar
Pliki konfiguracyjne PostgreSQL znajdują się w katalogu danych (\sphinxcode{\sphinxupquote{PGDATA}}).

\sphinxAtStartPar
\sphinxstylestrong{Główne pliki konfiguracyjne:}
\begin{itemize}
\item {} 
\sphinxAtStartPar
\sphinxcode{\sphinxupquote{postgresql.conf}} \textendash{} podstawowa konfiguracja serwera (port, pamięć, logi)

\item {} 
\sphinxAtStartPar
\sphinxcode{\sphinxupquote{pg\_hba.conf}} \textendash{} kontrola dostępu (kto i skąd może się łączyć)

\item {} 
\sphinxAtStartPar
\sphinxcode{\sphinxupquote{pg\_ident.conf}} \textendash{} mapowanie użytkowników systemowych na role PostgreSQL

\end{itemize}

\sphinxAtStartPar
\sphinxstylestrong{Domyślne lokalizacje:}


\begin{savenotes}\sphinxattablestart
\sphinxthistablewithglobalstyle
\centering
\sphinxcapstartof{table}
\sphinxthecaptionisattop
\sphinxcaption{Domyślne lokalizacje}\label{\detokenize{rozdzial_2/rozdzial_22/index:id6}}
\sphinxaftertopcaption
\begin{tabular}[t]{\X{30}{100}\X{70}{100}}
\sphinxtoprule
\sphinxstyletheadfamily 
\sphinxAtStartPar
System
&\sphinxstyletheadfamily 
\sphinxAtStartPar
Katalog danych
\\
\sphinxmidrule
\sphinxtableatstartofbodyhook
\sphinxAtStartPar
Linux (RPM/DEB)
&
\sphinxAtStartPar
\sphinxcode{\sphinxupquote{/var/lib/pgsql/data/}} lub \sphinxcode{\sphinxupquote{/var/lib/postgresql/*/main/}}
\\
\sphinxhline
\sphinxAtStartPar
Linux (kompilacja źródłowa)
&
\sphinxAtStartPar
\sphinxcode{\sphinxupquote{/usr/local/pgsql/data/}}
\\
\sphinxhline
\sphinxAtStartPar
Windows
&
\sphinxAtStartPar
\sphinxcode{\sphinxupquote{C:\textbackslash{}Program Files\textbackslash{}PostgreSQL\textbackslash{}\textless{}version\textgreater{}\textbackslash{}data\textbackslash{}}}
\\
\sphinxhline
\sphinxAtStartPar
Docker
&
\sphinxAtStartPar
\sphinxcode{\sphinxupquote{/var/lib/postgresql/data/}} (w kontenerze)
\\
\sphinxbottomrule
\end{tabular}
\sphinxtableafterendhook\par
\sphinxattableend\end{savenotes}

\sphinxAtStartPar
\sphinxstylestrong{Zmiana katalogu danych:}

\begin{sphinxVerbatim}[commandchars=\\\{\}]
\PYG{c+c1}{\PYGZsh{} Inicjalizacja nowego katalogu}
initdb\PYG{+w}{ }\PYGZhy{}D\PYG{+w}{ }/ścieżka/do/katalogu

\PYG{c+c1}{\PYGZsh{} Uruchomienie z niestandardowym PGDATA}
pg\PYGZus{}ctl\PYG{+w}{ }\PYGZhy{}D\PYG{+w}{ }/ścieżka/do/katalogu\PYG{+w}{ }start
\end{sphinxVerbatim}

\sphinxAtStartPar
\sphinxstylestrong{Struktura katalogu danych:}

\begin{sphinxVerbatim}[commandchars=\\\{\}]
\PYGZdl{}PGDATA/
├── postgresql.conf   \PYGZsh{} główny plik konfiguracyjny
├── pg\PYGZus{}hba.conf       \PYGZsh{} polityka dostępu
├── pg\PYGZus{}ident.conf     \PYGZsh{} mapowanie tożsamości
├── base/             \PYGZsh{} rzeczywiste bazy danych
├── global/           \PYGZsh{} tabele systemowe
├── log/              \PYGZsh{} logi serwera
└── pg\PYGZus{}wal/           \PYGZsh{} Write\PYGZhy{}Ahead Log (WAL)
\end{sphinxVerbatim}

\sphinxAtStartPar
\sphinxstylestrong{Sprawdzenie aktualnej lokalizacji:}

\begin{sphinxVerbatim}[commandchars=\\\{\}]
\PYG{k}{SHOW}\PYG{+w}{ }\PYG{n}{data\PYGZus{}directory}\PYG{p}{;}
\PYG{k}{SHOW}\PYG{+w}{ }\PYG{n}{config\PYGZus{}file}\PYG{p}{;}
\PYG{k}{SHOW}\PYG{+w}{ }\PYG{n}{hba\PYGZus{}file}\PYG{p}{;}
\end{sphinxVerbatim}


\subsection{Środowiska i profile}
\label{\detokenize{rozdzial_2/rozdzial_22/index:srodowiska-i-profile}}
\sphinxAtStartPar
Przełączanie środowisk PostgreSQL (dev/test/prod) realizuj przez:
\begin{itemize}
\item {} 
\sphinxAtStartPar
zmienne środowiskowe (np. \sphinxcode{\sphinxupquote{PGHOST}}, \sphinxcode{\sphinxupquote{PGPORT}}, \sphinxcode{\sphinxupquote{PGDATABASE}}, \sphinxcode{\sphinxupquote{PGUSER}}, \sphinxcode{\sphinxupquote{PGPASSWORD}}),

\item {} 
\sphinxAtStartPar
profile konfiguracyjne,

\item {} 
\sphinxAtStartPar
osobne pliki \sphinxcode{\sphinxupquote{.env}} dla każdego środowiska.

\end{itemize}


\subsection{Automatyzacja}
\label{\detokenize{rozdzial_2/rozdzial_22/index:automatyzacja}}
\sphinxAtStartPar
Zalecenia dla PostgreSQL:
\begin{itemize}
\item {} 
\sphinxAtStartPar
walidacja konfiguracji przy starcie za pomocą \sphinxcode{\sphinxupquote{pg\_isready}},

\item {} 
\sphinxAtStartPar
skrypty sprawdzające kompletność zmiennych środowiskowych,

\item {} 
\sphinxAtStartPar
integracja z CI/CD (np. GitHub Actions, GitLab CI).

\end{itemize}

\sphinxAtStartPar
Dokumentację generuj automatycznie (Sphinx).


\subsection{Przykłady}
\label{\detokenize{rozdzial_2/rozdzial_22/index:przyklady}}
\sphinxAtStartPar
\sphinxstylestrong{PostgreSQL \textendash{} zmienne środowiskowe (.env)}

\begin{sphinxVerbatim}[commandchars=\\\{\}]
PGHOST=localhost
PGPORT=5432
PGDATABASE=aplikacja
PGUSER=app\PYGZus{}user
PGPASSWORD=\PYGZdl{}\PYGZob{}DB\PYGZus{}PASSWORD\PYGZcb{}
\end{sphinxVerbatim}

\sphinxAtStartPar
\sphinxstylestrong{PostgreSQL \textendash{} URL połączenia}

\begin{sphinxVerbatim}[commandchars=\\\{\}]
DATABASE\PYGZus{}URL=postgresql://\PYGZdl{}\PYGZob{}PGUSER\PYGZcb{}:\PYGZdl{}\PYGZob{}PGPASSWORD\PYGZcb{}@\PYGZdl{}\PYGZob{}PGHOST\PYGZcb{}:\PYGZdl{}\PYGZob{}PGPORT\PYGZcb{}/\PYGZdl{}\PYGZob{}PGDATABASE\PYGZcb{}
\end{sphinxVerbatim}

\sphinxAtStartPar
\sphinxstylestrong{PostgreSQL \textendash{} sprawdzenie połączenia}

\begin{sphinxVerbatim}[commandchars=\\\{\}]
pg\PYGZus{}isready \PYGZhy{}h \PYGZdl{}\PYGZob{}PGHOST\PYGZcb{} \PYGZhy{}p \PYGZdl{}\PYGZob{}PGPORT\PYGZcb{} \PYGZhy{}U \PYGZdl{}\PYGZob{}PGUSER\PYGZcb{} \PYGZhy{}d \PYGZdl{}\PYGZob{}PGDATABASE\PYGZcb{}
\end{sphinxVerbatim}


\subsection{Najlepsze praktyki}
\label{\detokenize{rozdzial_2/rozdzial_22/index:najlepsze-praktyki}}\begin{enumerate}
\sphinxsetlistlabels{\arabic}{enumi}{enumii}{}{.}%
\item {} 
\sphinxAtStartPar
Nie przechowuj sekretów w repozytorium.

\item {} 
\sphinxAtStartPar
Stosuj \sphinxcode{\sphinxupquote{.env.example}} zamiast rzeczywistych danych.

\item {} 
\sphinxAtStartPar
Używaj standardowych zmiennych środowiskowych PostgreSQL (\sphinxcode{\sphinxupquote{PG*}}).

\item {} 
\sphinxAtStartPar
Waliduj konfigurację przed uruchomieniem (\sphinxcode{\sphinxupquote{pg\_isready}}).

\item {} 
\sphinxAtStartPar
W środowiskach produkcyjnych używaj połączeń SSL/TLS (\sphinxcode{\sphinxupquote{sslmode=require}}).

\end{enumerate}


\subsection{Podsumowanie}
\label{\detokenize{rozdzial_2/rozdzial_22/index:podsumowanie}}
\sphinxAtStartPar
Poprawna konfiguracja połączenia z PostgreSQL zwiększa bezpieczeństwo i przenośność aplikacji między środowiskami. Standardowe zmienne środowiskowe oraz narzędzia takie jak \sphinxcode{\sphinxupquote{pg\_isready}} ułatwiają automatyzację i walidację.


\subsection{Szczegółowe omówienie plików konfiguracyjnych}
\label{\detokenize{rozdzial_2/rozdzial_22/index:szczegolowe-omowienie-plikow-konfiguracyjnych}}
\sphinxAtStartPar
Po omówieniu podstawowych wymagań środowiskowych można przejść do szczegółowej konfiguracji PostgreSQL.

\sphinxAtStartPar
W kolejnych sekcjach przedstawiono najważniejsze parametry konfiguracyjne dostępne w plikach:
\begin{itemize}
\item {} 
\sphinxAtStartPar
\sphinxcode{\sphinxupquote{postgresql.conf}} \textendash{} konfiguracja działania serwera,

\item {} 
\sphinxAtStartPar
\sphinxcode{\sphinxupquote{pg\_hba.conf}} \textendash{} kontrola dostępu do bazy danych,

\item {} 
\sphinxAtStartPar
\sphinxcode{\sphinxupquote{pg\_ident.conf}} \textendash{} mapowanie użytkowników systemowych na role PostgreSQL.

\end{itemize}


\subsection{postgresql.conf}
\label{\detokenize{rozdzial_2/rozdzial_22/index:postgresql-conf}}
\sphinxAtStartPar
Plik \sphinxcode{\sphinxupquote{postgresql.conf}} jest głównym plikiem konfiguracyjnym klastra PostgreSQL.
Zawiera ustawienia wpływające na działanie serwera, wydajność, bezpieczeństwo, logowanie oraz replikację.

\sphinxAtStartPar
Lokalizację aktywnego pliku można sprawdzić poleceniem:

\begin{sphinxVerbatim}[commandchars=\\\{\}]
\PYG{k}{SHOW}\PYG{+w}{ }\PYG{n}{config\PYGZus{}file}\PYG{p}{;}
\end{sphinxVerbatim}


\subsubsection{Najważniejsze parametry}
\label{\detokenize{rozdzial_2/rozdzial_22/index:najwazniejsze-parametry}}

\paragraph{Połączenia sieciowe}
\label{\detokenize{rozdzial_2/rozdzial_22/index:polaczenia-sieciowe}}\begin{description}
\sphinxlineitem{\sphinxcode{\sphinxupquote{listen\_addresses}}}
\sphinxAtStartPar
Określa adresy IP, na których PostgreSQL nasłuchuje połączeń.

\sphinxAtStartPar
Najczęstsze wartości:
\begin{itemize}
\item {} 
\sphinxAtStartPar
\sphinxcode{\sphinxupquote{localhost}} \textendash{} tylko połączenia lokalne,

\item {} 
\sphinxAtStartPar
\sphinxcode{\sphinxupquote{*}} \textendash{} wszystkie interfejsy sieciowe,

\item {} 
\sphinxAtStartPar
lista adresów, np. \sphinxcode{\sphinxupquote{\textquotesingle{}192.168.1.10,localhost\textquotesingle{}}}.

\end{itemize}

\sphinxlineitem{\sphinxcode{\sphinxupquote{port}}}
\sphinxAtStartPar
Port TCP używany przez PostgreSQL.

\sphinxAtStartPar
Domyślna wartość:

\begin{sphinxVerbatim}[commandchars=\\\{\}]
\PYG{n+na}{port}\PYG{+w}{ }\PYG{o}{=}\PYG{+w}{ }\PYG{l+s}{5432}
\end{sphinxVerbatim}

\sphinxlineitem{\sphinxcode{\sphinxupquote{max\_connections}}}
\sphinxAtStartPar
Maksymalna liczba jednoczesnych połączeń z bazą danych.

\sphinxAtStartPar
Zbyt wysoka wartość może zwiększać zużycie pamięci.

\end{description}


\paragraph{Pamięć i wydajność}
\label{\detokenize{rozdzial_2/rozdzial_22/index:pamiec-i-wydajnosc}}\begin{description}
\sphinxlineitem{\sphinxcode{\sphinxupquote{shared\_buffers}}}
\sphinxAtStartPar
Ilość pamięci RAM używanej przez PostgreSQL do buforowania danych.

\sphinxAtStartPar
Typowe wartości:
\begin{itemize}
\item {} 
\sphinxAtStartPar
128MB \textendash{} środowiska testowe,

\item {} 
\sphinxAtStartPar
25\% pamięci RAM \textendash{} środowiska produkcyjne.

\end{itemize}

\sphinxlineitem{\sphinxcode{\sphinxupquote{work\_mem}}}
\sphinxAtStartPar
Ilość pamięci wykorzystywanej przez operacje sortowania i zapytania.

\sphinxAtStartPar
Parametr jest przydzielany dla pojedynczej operacji.

\sphinxlineitem{\sphinxcode{\sphinxupquote{maintenance\_work\_mem}}}
\sphinxAtStartPar
Pamięć używana podczas operacji administracyjnych, np.:
\begin{itemize}
\item {} 
\sphinxAtStartPar
\sphinxcode{\sphinxupquote{VACUUM}},

\item {} 
\sphinxAtStartPar
\sphinxcode{\sphinxupquote{CREATE INDEX}},

\item {} 
\sphinxAtStartPar
\sphinxcode{\sphinxupquote{ALTER TABLE}}.

\end{itemize}

\sphinxlineitem{\sphinxcode{\sphinxupquote{effective\_cache\_size}}}
\sphinxAtStartPar
Szacowany rozmiar pamięci podręcznej dostępnej dla PostgreSQL.

\sphinxAtStartPar
Parametr pomaga optymalizatorowi zapytań wybierać odpowiednie plany wykonania.

\end{description}


\paragraph{Logowanie}
\label{\detokenize{rozdzial_2/rozdzial_22/index:logowanie}}\begin{description}
\sphinxlineitem{\sphinxcode{\sphinxupquote{logging\_collector}}}
\sphinxAtStartPar
Włącza zapisywanie logów do plików.

\sphinxAtStartPar
Dostępne wartości:
\begin{itemize}
\item {} 
\sphinxAtStartPar
\sphinxcode{\sphinxupquote{on}},

\item {} 
\sphinxAtStartPar
\sphinxcode{\sphinxupquote{off}}.

\end{itemize}

\sphinxlineitem{\sphinxcode{\sphinxupquote{log\_directory}}}
\sphinxAtStartPar
Katalog przechowywania logów.

\sphinxlineitem{\sphinxcode{\sphinxupquote{log\_filename}}}
\sphinxAtStartPar
Wzorzec nazw plików logów.

\sphinxAtStartPar
Przykład:

\begin{sphinxVerbatim}[commandchars=\\\{\}]
\PYG{n+na}{log\PYGZus{}filename}\PYG{+w}{ }\PYG{o}{=}\PYG{+w}{ }\PYG{l+s}{\PYGZsq{}postgresql\PYGZhy{}\PYGZpc{}Y\PYGZhy{}\PYGZpc{}m\PYGZhy{}\PYGZpc{}d.log\PYGZsq{}}
\end{sphinxVerbatim}

\sphinxlineitem{\sphinxcode{\sphinxupquote{log\_min\_duration\_statement}}}
\sphinxAtStartPar
Określa minimalny czas wykonania zapytania (w ms), po którym zostanie ono zapisane w logach.

\sphinxAtStartPar
Przykład:

\begin{sphinxVerbatim}[commandchars=\\\{\}]
\PYG{n+na}{log\PYGZus{}min\PYGZus{}duration\PYGZus{}statement}\PYG{+w}{ }\PYG{o}{=}\PYG{+w}{ }\PYG{l+s}{1000}
\end{sphinxVerbatim}

\sphinxAtStartPar
Wartość \sphinxcode{\sphinxupquote{1000}} oznacza logowanie zapytań trwających dłużej niż 1 sekundę.

\end{description}


\paragraph{Bezpieczeństwo}
\label{\detokenize{rozdzial_2/rozdzial_22/index:bezpieczenstwo}}\begin{description}
\sphinxlineitem{\sphinxcode{\sphinxupquote{password\_encryption}}}
\sphinxAtStartPar
Algorytm szyfrowania haseł użytkowników.

\sphinxAtStartPar
Zalecana wartość:

\begin{sphinxVerbatim}[commandchars=\\\{\}]
\PYG{n+na}{password\PYGZus{}encryption}\PYG{+w}{ }\PYG{o}{=}\PYG{+w}{ }\PYG{l+s}{scram\PYGZhy{}sha\PYGZhy{}256}
\end{sphinxVerbatim}

\sphinxlineitem{\sphinxcode{\sphinxupquote{ssl}}}
\sphinxAtStartPar
Włącza obsługę połączeń SSL/TLS.

\sphinxAtStartPar
Dostępne wartości:
\begin{itemize}
\item {} 
\sphinxAtStartPar
\sphinxcode{\sphinxupquote{on}},

\item {} 
\sphinxAtStartPar
\sphinxcode{\sphinxupquote{off}}.

\end{itemize}

\sphinxlineitem{\sphinxcode{\sphinxupquote{ssl\_cert\_file}} / \sphinxcode{\sphinxupquote{ssl\_key\_file}}}
\sphinxAtStartPar
Ścieżki do certyfikatu i klucza prywatnego SSL.

\end{description}


\paragraph{Replikacja i WAL}
\label{\detokenize{rozdzial_2/rozdzial_22/index:replikacja-i-wal}}\begin{description}
\sphinxlineitem{\sphinxcode{\sphinxupquote{wal\_level}}}
\sphinxAtStartPar
Określa poziom szczegółowości Write\sphinxhyphen{}Ahead Log (WAL).

\sphinxAtStartPar
Dostępne wartości:
\begin{itemize}
\item {} 
\sphinxAtStartPar
\sphinxcode{\sphinxupquote{minimal}},

\item {} 
\sphinxAtStartPar
\sphinxcode{\sphinxupquote{replica}},

\item {} 
\sphinxAtStartPar
\sphinxcode{\sphinxupquote{logical}}.

\end{itemize}

\sphinxlineitem{\sphinxcode{\sphinxupquote{max\_wal\_size}}}
\sphinxAtStartPar
Maksymalny rozmiar plików WAL przed wykonaniem checkpointu.

\sphinxlineitem{\sphinxcode{\sphinxupquote{archive\_mode}}}
\sphinxAtStartPar
Włącza archiwizację WAL.

\sphinxlineitem{\sphinxcode{\sphinxupquote{archive\_command}}}
\sphinxAtStartPar
Polecenie wykonywane podczas archiwizacji plików WAL.

\end{description}


\subsubsection{Przykładowa konfiguracja}
\label{\detokenize{rozdzial_2/rozdzial_22/index:przykladowa-konfiguracja}}
\begin{sphinxVerbatim}[commandchars=\\\{\}]
\PYG{n+na}{listen\PYGZus{}addresses}\PYG{+w}{ }\PYG{o}{=}\PYG{+w}{ }\PYG{l+s}{\PYGZsq{}*\PYGZsq{}}
\PYG{n+na}{port}\PYG{+w}{ }\PYG{o}{=}\PYG{+w}{ }\PYG{l+s}{5432}
\PYG{n+na}{max\PYGZus{}connections}\PYG{+w}{ }\PYG{o}{=}\PYG{+w}{ }\PYG{l+s}{200}

\PYG{n+na}{shared\PYGZus{}buffers}\PYG{+w}{ }\PYG{o}{=}\PYG{+w}{ }\PYG{l+s}{1GB}
\PYG{n+na}{work\PYGZus{}mem}\PYG{+w}{ }\PYG{o}{=}\PYG{+w}{ }\PYG{l+s}{16MB}
\PYG{n+na}{maintenance\PYGZus{}work\PYGZus{}mem}\PYG{+w}{ }\PYG{o}{=}\PYG{+w}{ }\PYG{l+s}{256MB}

\PYG{n+na}{logging\PYGZus{}collector}\PYG{+w}{ }\PYG{o}{=}\PYG{+w}{ }\PYG{l+s}{on}
\PYG{n+na}{log\PYGZus{}directory}\PYG{+w}{ }\PYG{o}{=}\PYG{+w}{ }\PYG{l+s}{\PYGZsq{}log\PYGZsq{}}

\PYG{n+na}{ssl}\PYG{+w}{ }\PYG{o}{=}\PYG{+w}{ }\PYG{l+s}{on}
\PYG{n+na}{password\PYGZus{}encryption}\PYG{+w}{ }\PYG{o}{=}\PYG{+w}{ }\PYG{l+s}{scram\PYGZhy{}sha\PYGZhy{}256}

\PYG{n+na}{wal\PYGZus{}level}\PYG{+w}{ }\PYG{o}{=}\PYG{+w}{ }\PYG{l+s}{replica}
\PYG{n+na}{max\PYGZus{}wal\PYGZus{}size}\PYG{+w}{ }\PYG{o}{=}\PYG{+w}{ }\PYG{l+s}{2GB}
\end{sphinxVerbatim}


\subsubsection{Weryfikacja konfiguracji}
\label{\detokenize{rozdzial_2/rozdzial_22/index:weryfikacja-konfiguracji}}
\sphinxAtStartPar
Po zmianie parametrów konfigurację można przeładować bez restartu serwera:

\begin{sphinxVerbatim}[commandchars=\\\{\}]
SELECT\PYG{+w}{ }pg\PYGZus{}reload\PYGZus{}conf\PYG{o}{(}\PYG{o}{)}\PYG{p}{;}
\end{sphinxVerbatim}

\sphinxAtStartPar
Niektóre parametry wymagają pełnego restartu PostgreSQL.

\sphinxAtStartPar
Aktualne wartości parametrów można sprawdzić poleceniem:

\begin{sphinxVerbatim}[commandchars=\\\{\}]
\PYG{k}{SHOW}\PYG{+w}{ }\PYG{n}{shared\PYGZus{}buffers}\PYG{p}{;}
\PYG{k}{SHOW}\PYG{+w}{ }\PYG{n}{wal\PYGZus{}level}\PYG{p}{;}
\PYG{k}{SHOW}\PYG{+w}{ }\PYG{n}{max\PYGZus{}connections}\PYG{p}{;}
\end{sphinxVerbatim}


\subsection{pg\_hba.conf}
\label{\detokenize{rozdzial_2/rozdzial_22/index:pg-hba-conf}}
\sphinxAtStartPar
Plik \sphinxcode{\sphinxupquote{pg\_hba.conf}} (Host\sphinxhyphen{}Based Authentication) odpowiada za kontrolę dostępu do serwera PostgreSQL.

\sphinxAtStartPar
Definiuje:
\begin{itemize}
\item {} 
\sphinxAtStartPar
kto może połączyć się z bazą danych,

\item {} 
\sphinxAtStartPar
z jakiego adresu IP,

\item {} 
\sphinxAtStartPar
do jakiej bazy danych,

\item {} 
\sphinxAtStartPar
przy użyciu jakiej metody uwierzytelniania.

\end{itemize}

\sphinxAtStartPar
Lokalizację aktywnego pliku można sprawdzić poleceniem:

\begin{sphinxVerbatim}[commandchars=\\\{\}]
\PYG{k}{SHOW}\PYG{+w}{ }\PYG{n}{hba\PYGZus{}file}\PYG{p}{;}
\end{sphinxVerbatim}


\subsubsection{Składnia wpisu}
\label{\detokenize{rozdzial_2/rozdzial_22/index:skladnia-wpisu}}
\sphinxAtStartPar
Każdy wpis w \sphinxcode{\sphinxupquote{pg\_hba.conf}} ma następującą strukturę:

\begin{sphinxVerbatim}[commandchars=\\\{\}]
\PYG{n}{TYPE}    \PYG{n}{DATABASE}    \PYG{n}{USER}    \PYG{n}{ADDRESS}    \PYG{n}{METHOD}
\end{sphinxVerbatim}

\sphinxAtStartPar
Znaczenie pól:
\begin{description}
\sphinxlineitem{\sphinxcode{\sphinxupquote{TYPE}}}
\sphinxAtStartPar
Typ połączenia.

\sphinxAtStartPar
Dostępne wartości:
\begin{itemize}
\item {} 
\sphinxAtStartPar
\sphinxcode{\sphinxupquote{local}} \textendash{} połączenia lokalne (Unix socket),

\item {} 
\sphinxAtStartPar
\sphinxcode{\sphinxupquote{host}} \textendash{} połączenia TCP/IP,

\item {} 
\sphinxAtStartPar
\sphinxcode{\sphinxupquote{hostssl}} \textendash{} połączenia TCP/IP z SSL,

\item {} 
\sphinxAtStartPar
\sphinxcode{\sphinxupquote{hostnossl}} \textendash{} połączenia TCP/IP bez SSL.

\end{itemize}

\sphinxlineitem{\sphinxcode{\sphinxupquote{DATABASE}}}
\sphinxAtStartPar
Nazwa bazy danych, do której użytkownik może się połączyć.

\sphinxAtStartPar
Możliwe wartości:
\begin{itemize}
\item {} 
\sphinxAtStartPar
konkretna nazwa bazy,

\item {} 
\sphinxAtStartPar
\sphinxcode{\sphinxupquote{all}} \textendash{} wszystkie bazy,

\item {} 
\sphinxAtStartPar
\sphinxcode{\sphinxupquote{sameuser}} \textendash{} baza o nazwie zgodnej z użytkownikiem.

\end{itemize}

\sphinxlineitem{\sphinxcode{\sphinxupquote{USER}}}
\sphinxAtStartPar
Użytkownik lub rola PostgreSQL.

\sphinxlineitem{\sphinxcode{\sphinxupquote{ADDRESS}}}
\sphinxAtStartPar
Adres IP lub zakres sieci w notacji CIDR.

\sphinxAtStartPar
Przykłady:
\begin{itemize}
\item {} 
\sphinxAtStartPar
\sphinxcode{\sphinxupquote{127.0.0.1/32}},

\item {} 
\sphinxAtStartPar
\sphinxcode{\sphinxupquote{192.168.1.0/24}},

\item {} 
\sphinxAtStartPar
\sphinxcode{\sphinxupquote{0.0.0.0/0}} \textendash{} wszystkie adresy.

\end{itemize}

\sphinxlineitem{\sphinxcode{\sphinxupquote{METHOD}}}
\sphinxAtStartPar
Metoda uwierzytelniania użytkownika.

\end{description}


\subsubsection{Najczęściej używane metody uwierzytelniania}
\label{\detokenize{rozdzial_2/rozdzial_22/index:najczesciej-uzywane-metody-uwierzytelniania}}\begin{description}
\sphinxlineitem{\sphinxcode{\sphinxupquote{trust}}}
\sphinxAtStartPar
Zezwala na połączenie bez uwierzytelniania.

\sphinxAtStartPar
Zalecane wyłącznie w środowiskach testowych.

\sphinxlineitem{\sphinxcode{\sphinxupquote{reject}}}
\sphinxAtStartPar
Odrzuca połączenie.

\sphinxlineitem{\sphinxcode{\sphinxupquote{md5}}}
\sphinxAtStartPar
Uwierzytelnianie przy użyciu hasła MD5.

\sphinxAtStartPar
Metoda starsza i mniej bezpieczna niż SCRAM.

\sphinxlineitem{\sphinxcode{\sphinxupquote{scram\sphinxhyphen{}sha\sphinxhyphen{}256}}}
\sphinxAtStartPar
Nowoczesna metoda uwierzytelniania oparta na SCRAM.

\sphinxAtStartPar
Zalecana dla nowych instalacji PostgreSQL.

\sphinxlineitem{\sphinxcode{\sphinxupquote{peer}}}
\sphinxAtStartPar
Uwierzytelnianie na podstawie użytkownika systemowego.

\sphinxAtStartPar
Najczęściej używane dla lokalnych połączeń Linux/Unix.

\sphinxlineitem{\sphinxcode{\sphinxupquote{ident}}}
\sphinxAtStartPar
Mapowanie użytkownika systemowego przy użyciu \sphinxcode{\sphinxupquote{pg\_ident.conf}}.

\end{description}


\subsubsection{Przykładowa konfiguracja}
\label{\detokenize{rozdzial_2/rozdzial_22/index:id1}}
\begin{sphinxVerbatim}[commandchars=\\\{\}]
\PYG{c+c1}{\PYGZsh{} Połączenia lokalne}
\PYG{n+na}{local   all             postgres                        peer}

\PYG{c+c1}{\PYGZsh{} Połączenia localhost IPv4}
\PYG{n+na}{host    all             all         127.0.0.1/32        scram\PYGZhy{}sha\PYGZhy{}256}

\PYG{c+c1}{\PYGZsh{} Połączenia localhost IPv6}
\PYG{n+na}{host    all             all}\PYG{+w}{         }\PYG{o}{:}\PYG{l+s}{:1/128             scram\PYGZhy{}sha\PYGZhy{}256}

\PYG{c+c1}{\PYGZsh{} Sieć lokalna}
\PYG{n+na}{host    aplikacja       app\PYGZus{}user    192.168.1.0/24      scram\PYGZhy{}sha\PYGZhy{}256}
\end{sphinxVerbatim}


\subsubsection{Najlepsze praktyki}
\label{\detokenize{rozdzial_2/rozdzial_22/index:id2}}\begin{enumerate}
\sphinxsetlistlabels{\arabic}{enumi}{enumii}{}{.}%
\item {} 
\sphinxAtStartPar
Używaj \sphinxcode{\sphinxupquote{scram\sphinxhyphen{}sha\sphinxhyphen{}256}} zamiast \sphinxcode{\sphinxupquote{md5}}.

\item {} 
\sphinxAtStartPar
Ograniczaj dostęp do konkretnych adresów IP.

\item {} 
\sphinxAtStartPar
Nie używaj \sphinxcode{\sphinxupquote{trust}} w środowiskach produkcyjnych.

\item {} 
\sphinxAtStartPar
Stosuj \sphinxcode{\sphinxupquote{hostssl}} dla połączeń zdalnych.

\item {} 
\sphinxAtStartPar
Po zmianie konfiguracji wykonuj reload konfiguracji:

\end{enumerate}

\begin{sphinxVerbatim}[commandchars=\\\{\}]
\PYG{k}{SELECT}\PYG{+w}{ }\PYG{n}{pg\PYGZus{}reload\PYGZus{}conf}\PYG{p}{(}\PYG{p}{)}\PYG{p}{;}
\end{sphinxVerbatim}


\subsection{pg\_ident.conf}
\label{\detokenize{rozdzial_2/rozdzial_22/index:pg-ident-conf}}
\sphinxAtStartPar
Plik \sphinxcode{\sphinxupquote{pg\_ident.conf}} umożliwia mapowanie użytkowników systemowych na role PostgreSQL.

\sphinxAtStartPar
Najczęściej używany jest razem z metodami uwierzytelniania:
\begin{itemize}
\item {} 
\sphinxAtStartPar
\sphinxcode{\sphinxupquote{peer}},

\item {} 
\sphinxAtStartPar
\sphinxcode{\sphinxupquote{ident}}.

\end{itemize}

\sphinxAtStartPar
Lokalizację aktywnego pliku można sprawdzić poleceniem:

\begin{sphinxVerbatim}[commandchars=\\\{\}]
\PYG{k}{SHOW}\PYG{+w}{ }\PYG{n}{ident\PYGZus{}file}\PYG{p}{;}
\end{sphinxVerbatim}


\subsubsection{Zastosowanie}
\label{\detokenize{rozdzial_2/rozdzial_22/index:zastosowanie}}
\sphinxAtStartPar
Mechanizm mapowania pozwala określić, który użytkownik systemowy może zostać powiązany z konkretną rolą PostgreSQL.

\sphinxAtStartPar
Jest to szczególnie przydatne:
\begin{itemize}
\item {} 
\sphinxAtStartPar
w środowiskach Linux/Unix,

\item {} 
\sphinxAtStartPar
przy automatyzacji administracyjnej,

\item {} 
\sphinxAtStartPar
dla lokalnych połączeń bez użycia haseł.

\end{itemize}


\subsubsection{Składnia wpisu}
\label{\detokenize{rozdzial_2/rozdzial_22/index:id3}}
\begin{sphinxVerbatim}[commandchars=\\\{\}]
\PYG{n}{MAPNAME}    \PYG{n}{SYSTEM}\PYG{o}{\PYGZhy{}}\PYG{n}{USERNAME}    \PYG{n}{PG}\PYG{o}{\PYGZhy{}}\PYG{n}{USERNAME}
\end{sphinxVerbatim}

\sphinxAtStartPar
Znaczenie pól:
\begin{description}
\sphinxlineitem{\sphinxcode{\sphinxupquote{MAPNAME}}}
\sphinxAtStartPar
Nazwa mapowania wykorzystywana w \sphinxcode{\sphinxupquote{pg\_hba.conf}}.

\sphinxlineitem{\sphinxcode{\sphinxupquote{SYSTEM\sphinxhyphen{}USERNAME}}}
\sphinxAtStartPar
Użytkownik systemu operacyjnego.

\sphinxlineitem{\sphinxcode{\sphinxupquote{PG\sphinxhyphen{}USERNAME}}}
\sphinxAtStartPar
Rola PostgreSQL, na którą użytkownik ma zostać zmapowany.

\end{description}


\subsubsection{Przykładowa konfiguracja}
\label{\detokenize{rozdzial_2/rozdzial_22/index:id4}}
\begin{sphinxVerbatim}[commandchars=\\\{\}]
\PYG{c+c1}{\PYGZsh{} MAPNAME    SYSTEM\PYGZhy{}USERNAME    PG\PYGZhy{}USERNAME}
\PYG{n+na}{admin\PYGZus{}map    postgres           postgres}
\PYG{n+na}{admin\PYGZus{}map    deploy             db\PYGZus{}admin}
\PYG{n+na}{app\PYGZus{}map      appuser            aplikacja}
\end{sphinxVerbatim}

\sphinxAtStartPar
Przykład użycia w \sphinxcode{\sphinxupquote{pg\_hba.conf}}:

\begin{sphinxVerbatim}[commandchars=\\\{\}]
\PYG{n+na}{local   all     all     peer map}\PYG{o}{=}\PYG{l+s}{admin\PYGZus{}map}
\end{sphinxVerbatim}

\sphinxAtStartPar
W powyższym przykładzie użytkownicy systemowi zdefiniowani w \sphinxcode{\sphinxupquote{admin\_map}} mogą logować się jako odpowiadające im role PostgreSQL.


\subsubsection{Najlepsze praktyki}
\label{\detokenize{rozdzial_2/rozdzial_22/index:id5}}\begin{enumerate}
\sphinxsetlistlabels{\arabic}{enumi}{enumii}{}{.}%
\item {} 
\sphinxAtStartPar
Ograniczaj mapowania wyłącznie do wymaganych użytkowników.

\item {} 
\sphinxAtStartPar
Nie mapuj użytkowników systemowych na role superuser bez uzasadnienia.

\item {} 
\sphinxAtStartPar
Stosuj osobne role administracyjne i aplikacyjne.

\item {} 
\sphinxAtStartPar
Dokumentuj wszystkie mapowania użytkowników.

\item {} 
\sphinxAtStartPar
Regularnie przeglądaj konfigurację dostępu.

\end{enumerate}

\begin{sphinxadmonition}{note}{Opracowanie}

\sphinxAtStartPar
\sphinxstylestrong{Autor:} Michał Kraus
\sphinxstylestrong{Przedmiot:} Bazy danych
\sphinxstylestrong{Data:} maj 2026
\end{sphinxadmonition}

\sphinxstepscope


\section{Kontrola i konserwacja}
\label{\detokenize{rozdzial_2/rozdzial_23/index:kontrola-i-konserwacja}}\label{\detokenize{rozdzial_2/rozdzial_23/index::doc}}\begin{quote}\begin{description}
\sphinxlineitem{Autorzy}\begin{enumerate}
\sphinxsetlistlabels{\arabic}{enumi}{enumii}{}{.}%
\item {} 
\sphinxAtStartPar
Paweł Łoćwin

\item {} 
\sphinxAtStartPar
Paweł Łosowski

\end{enumerate}

\end{description}\end{quote}


\subsection{Wprowadzenie}
\label{\detokenize{rozdzial_2/rozdzial_23/index:wprowadzenie}}
\sphinxAtStartPar
Współczesne systemy relacyjnych baz danych to wysoce skomplikowane środowiska, które wymagają ciągłego i proaktywnego nadzoru, aby zagwarantować wysoką dostępność, niezawodność oraz wydajność. Administracja takimi systemami stanowi złożone zagadnienie wymagające zarówno wiedzy teoretycznej, jak i praktycznego doświadczenia w zarządzaniu infrastrukturą informatyczną.


\subsection{Zarządzanie przestrzenią i mechanizm współbieżności}
\label{\detokenize{rozdzial_2/rozdzial_23/index:zarzadzanie-przestrzenia-i-mechanizm-wspolbieznosci}}
\sphinxAtStartPar
PostgreSQL opiera swoje działanie na zaawansowanej architekturze określanej jako Multi\sphinxhyphen{}Version Concurrency Control (MVCC). Technologia ta pozwala na równoległy odczyt i zapis danych bez wzajemnego blokowania się transakcji, co stanowi fundamentalną zaletę w środowiskach wieloużytkownikowych o wysokim natężeniu operacji.

\sphinxAtStartPar
Aby zapobiec zjawisku drastycznego puchnięcia tabel oraz degradacji wydajności operacji wejścia i wyjścia, konieczna jest regularna konserwacja na poziomie nośnika danych:
\begin{itemize}
\item {} 
\sphinxAtStartPar
\sphinxstylestrong{VACUUM:} Podstawowy proces konserwacyjny odzyskujący miejsce po martwych krotkach. Oznacza on wolne obszary jako dostępne do ponownego zapisu przez nowe instrukcje, zapobiegając niekontrolowanemu wzrostowi rozmiarów tabel na dysku.

\item {} 
\sphinxAtStartPar
\sphinxstylestrong{VACUUM FULL:} Znacznie bardziej inwazyjna wersja operacji, która fizycznie przebudowuje strukturę całej tabeli i trwale zwraca wolne miejsce do systemu operacyjnego. Operacja ta wymaga jednak zablokowania całej tabeli i powinna być wykonywana w godzinach niskiego obciążenia.

\item {} 
\sphinxAtStartPar
\sphinxstylestrong{AUTOVACUUM:} W nowoczesnych środowiskach produkcyjnych utrzymanie czystości dyskowej powierza się wbudowanemu procesowi pracującemu w tle. Monitoruje on na bieżąco statystyki zmian i automatycznie wyzwala operacje czyszczenia, eliminując potrzebę ręcznej interwencji administratora.

\end{itemize}

\sphinxAtStartPar
Zaniedbanie procesu czyszczenia w systemach o wysokiej rotacji danych może doprowadzić nie tylko do spadku wydajności, ale w skrajnych przypadkach do wyczerpania identyfikatorów transakcji, co prowadzi do całkowitego zamrożenia bazy.

\begin{sphinxVerbatim}[commandchars=\\\{\}]
\PYG{c+c1}{\PYGZhy{}\PYGZhy{} Przykładowe wymuszenie manualnej konserwacji z analizą statystyk dla tabeli}
\PYG{k}{VACUUM}\PYG{+w}{ }\PYG{p}{(}\PYG{k}{VERBOSE}\PYG{p}{,}\PYG{+w}{ }\PYG{k}{ANALYZE}\PYG{p}{)}\PYG{+w}{ }\PYG{n}{Wypozyczenia}\PYG{p}{;}
\end{sphinxVerbatim}


\subsection{Optymalizacja i statystyki planisty zapytań}
\label{\detokenize{rozdzial_2/rozdzial_23/index:optymalizacja-i-statystyki-planisty-zapytan}}
\sphinxAtStartPar
Kolejnym filarem utrzymania sprawności systemu jest kontrola nad optymalizatorem, zwanym planistą zapytań. Silnik bazy danych przed wykonaniem jakiejkolwiek skomplikowanej kwerendy analizuje dostępne drogi realizacji i wybiera najefektywniejszą z punktu widzenia kosztów obliczeniowych.

\sphinxAtStartPar
Aby planista mógł podejmować optymalne decyzje, musi opierać się na wysoce dokładnych i aktualnych statystykach dotyczących rozkładu danych. Obejmują one między innymi histogramy wartości w kolumnach, liczby odrębnych wartości (distinct count), oraz rozmiary tabel i indeksów.
\begin{itemize}
\item {} 
\sphinxAtStartPar
Instrukcja \sphinxstylestrong{ANALYZE} służy do natychmiastowego odświeżenia tych metadanych. Baza pobiera losową próbkę wierszy i na jej podstawie aktualizuje wewnętrzne tabele systemowe, co pozwala uniknąć błędnych oszacowań planu wykonania.

\item {} 
\sphinxAtStartPar
Bieżące profilowanie wydajności realizowane jest natomiast przy pomocy instrukcji \sphinxstylestrong{EXPLAIN ANALYZE}. Narzędzie to uruchamia zapytanie, po czym zwraca nie tylko wynik, ale również szczegółowy opis planu wykonania wraz z rzeczywistymi czasy i liczbami wierszy na każdym kroku.

\end{itemize}

\begin{sphinxVerbatim}[commandchars=\\\{\}]
\PYG{c+c1}{\PYGZhy{}\PYGZhy{} Kontrola planu wykonania zapytania z rzeczywistym pomiarem czasu}
\PYG{k}{EXPLAIN}\PYG{+w}{ }\PYG{k}{ANALYZE}
\PYG{k}{SELECT}\PYG{+w}{ }\PYG{k}{c}\PYG{p}{.}\PYG{n}{Imie}\PYG{p}{,}\PYG{+w}{ }\PYG{k}{c}\PYG{p}{.}\PYG{n}{Nazwisko}\PYG{p}{,}\PYG{+w}{ }\PYG{n}{k}\PYG{p}{.}\PYG{n}{Tytul}
\PYG{k}{FROM}\PYG{+w}{ }\PYG{n}{Czytelnicy}\PYG{+w}{ }\PYG{k}{c}
\PYG{k}{JOIN}\PYG{+w}{ }\PYG{n}{Wypozyczenia}\PYG{+w}{ }\PYG{n}{w}\PYG{+w}{ }\PYG{k}{ON}\PYG{+w}{ }\PYG{k}{c}\PYG{p}{.}\PYG{n}{ID\PYGZus{}Czytelnika}\PYG{+w}{ }\PYG{o}{=}\PYG{+w}{ }\PYG{n}{w}\PYG{p}{.}\PYG{n}{ID\PYGZus{}Czytelnika}
\PYG{k}{JOIN}\PYG{+w}{ }\PYG{n}{Ksiazki}\PYG{+w}{ }\PYG{n}{k}\PYG{+w}{ }\PYG{k}{ON}\PYG{+w}{ }\PYG{n}{w}\PYG{p}{.}\PYG{n}{ID\PYGZus{}Ksiazki}\PYG{+w}{ }\PYG{o}{=}\PYG{+w}{ }\PYG{n}{k}\PYG{p}{.}\PYG{n}{ID\PYGZus{}Ksiazki}
\PYG{k}{WHERE}\PYG{+w}{ }\PYG{n}{w}\PYG{p}{.}\PYG{n}{Data\PYGZus{}Zwrotu}\PYG{+w}{ }\PYG{k}{IS}\PYG{+w}{ }\PYG{k}{NULL}\PYG{p}{;}
\end{sphinxVerbatim}


\subsection{Bezpieczeństwo i rygorystyczna kontrola dostępu}
\label{\detokenize{rozdzial_2/rozdzial_23/index:bezpieczenstwo-i-rygorystyczna-kontrola-dostepu}}
\sphinxAtStartPar
Bezpieczeństwo danych stanowi absolutny priorytet każdej administracji systemami informatycznymi. Nowoczesne podejście do ochrony baz danych całkowicie odchodzi od wykorzystywania globalnych kont administratora na rzecz zaawansowanego systemu uprawnień opartego na rolach (RBAC \textendash{} Role\sphinxhyphen{}Based Access Control).

\sphinxAtStartPar
System zabezpieczeń środowiska PostgreSQL pozwala na niezwykle elastyczną gradację uprawnień. Definiuje się dedykowane role systemowe przypisane do konkretnych mikrousług lub modułów aplikacji, z precyzyjnie określonymi prawami dostępu do pojedynczych schematów, tabel, a nawet kolumn.

\sphinxAtStartPar
Kluczowym konceptem jest wdrożenie zasady najmniejszych uprawnień (Principle of Least Privilege). Rola wykorzystywana przez aplikację powinna dysponować prawami pozwalającymi wyłącznie na modyfikację danych niezbędnych do jej funkcjonowania, bez dostępu do danych wrażliwych czy możliwości zmiany struktury bazy.

\begin{sphinxVerbatim}[commandchars=\\\{\}]
\PYG{c+c1}{\PYGZhy{}\PYGZhy{} Przykład implementacji bezpiecznej polityki kontroli dostępu}
\PYG{k}{CREATE}\PYG{+w}{ }\PYG{k}{ROLE}\PYG{+w}{ }\PYG{n}{app\PYGZus{}user}\PYG{+w}{ }\PYG{k}{WITH}\PYG{+w}{ }\PYG{n}{LOGIN}\PYG{+w}{ }\PYG{n}{PASSWORD}\PYG{+w}{ }\PYG{l+s+s1}{\PYGZsq{}TrudneHaslo123\PYGZsq{}}\PYG{p}{;}
\PYG{k}{GRANT}\PYG{+w}{ }\PYG{k}{CONNECT}\PYG{+w}{ }\PYG{k}{ON}\PYG{+w}{ }\PYG{k}{DATABASE}\PYG{+w}{ }\PYG{n}{biblioteka\PYGZus{}db}\PYG{+w}{ }\PYG{k}{TO}\PYG{+w}{ }\PYG{n}{app\PYGZus{}user}\PYG{p}{;}
\PYG{k}{GRANT}\PYG{+w}{ }\PYG{k}{USAGE}\PYG{+w}{ }\PYG{k}{ON}\PYG{+w}{ }\PYG{k}{SCHEMA}\PYG{+w}{ }\PYG{k}{public}\PYG{+w}{ }\PYG{k}{TO}\PYG{+w}{ }\PYG{n}{app\PYGZus{}user}\PYG{p}{;}
\PYG{k}{GRANT}\PYG{+w}{ }\PYG{k}{SELECT}\PYG{p}{,}\PYG{+w}{ }\PYG{k}{INSERT}\PYG{+w}{ }\PYG{k}{ON}\PYG{+w}{ }\PYG{n}{Wypozyczenia}\PYG{+w}{ }\PYG{k}{TO}\PYG{+w}{ }\PYG{n}{app\PYGZus{}user}\PYG{p}{;}
\end{sphinxVerbatim}


\subsection{Strategie zabezpieczania przed utratą danych}
\label{\detokenize{rozdzial_2/rozdzial_23/index:strategie-zabezpieczania-przed-utrata-danych}}
\sphinxAtStartPar
Nawet najlepiej zoptymalizowany i zabezpieczony system informatyczny jest narażony na błędy ludzkie lub awarie infrastruktury sprzętowej. Dlatego odpowiednia polityka wykonywania kopii zapasowych stanowi ostateczną linię obrony i jest absolutnie niezbędna dla każdej produkcyjnej instancji bazy danych.

\sphinxAtStartPar
Kopie logiczne polegają na ekstrakcji danych z bazy do postaci czystego skryptu języka SQL lub dedykowanego formatu archiwalnego. Narzędzia realizujące to zadanie gwarantują przenośność danych między platformami i wersjami silnika, jednak proces przywracania jest czasochłonny i może nie przywrócić wszystkich zmian z ostatnich minut przed awarią.

\sphinxAtStartPar
Dlatego w krytycznych środowiskach produkcyjnych stosuje się kopie fizyczne połączone z archiwizacją dzienników wyprzedzających. Dzienniki te rejestrują absolutnie każdą zmianę bajtu na poziomie magazynu, co umożliwia odtworzenie bazy do dowolnego punktu w przeszłości z dokładnością do sekundy.


\subsection{Podsumowanie}
\label{\detokenize{rozdzial_2/rozdzial_23/index:podsumowanie}}
\sphinxAtStartPar
Prawidłowa kontrola i konserwacja środowiska bazodanowego to wielowymiarowy proces, który nie ma punktu końcowego, lecz trwa przez cały cykl życia oprogramowania. Złożoność nowoczesnych systemów relacyjnych wymaga holistycznego podejścia łączącego automatyzację, monitoring i ciągłe doskonalenie.

\sphinxAtStartPar
Samo zaprojektowanie zgodnego ze sztuką schematu relacyjnego jest zaledwie początkiem pracy. Stabilność aplikacji zależy w równej mierze od rygorystycznego zarządzania fizyczną strukturą danych, optymalizacji wykonywanych operacji oraz wdrażania wielowarstwowych strategii bezpieczeństwa i ochrony przed awarią.

\sphinxAtStartPar
Holistyczne spojrzenie na administrację PostgreSQL, obejmujące zarówno automatyzację procesów porządkujących pamięć, jak i ciągłe monitorowanie planów wykonania zapytań, pozwala na budowanie systemów, które są nie tylko funkcjonalne, ale także odpornie i niezawodne w najbardziej wymagających warunkach produkcyjnych.

\sphinxstepscope

\sphinxAtStartPar
Tego typu jakis tekst

\sphinxstepscope


\section{Partycjonowanie danych}
\label{\detokenize{rozdzial_2/rozdzial_26/index:partycjonowanie-danych}}\label{\detokenize{rozdzial_2/rozdzial_26/index::doc}}
\sphinxstepscope


\section{Bezpieczeństwo Baz Danych}
\label{\detokenize{rozdzial_2/rozdzial_27/index:bezpieczenstwo-baz-danych}}\label{\detokenize{rozdzial_2/rozdzial_27/index::doc}}
\sphinxAtStartPar
Dział 7 przeglądu literaturowego poświęcony architekturze bezpieczeństwa, mechanizmom ochrony oraz analizie podatności w nowoczesnych systemach zarządzania bazami danych.

\sphinxstepscope


\subsection{1. Model bezpieczeństwa CIA i ochrona infrastruktury}
\label{\detokenize{rozdzial_2/rozdzial_27/rozdzial_1:model-bezpieczenstwa-cia-i-ochrona-infrastruktury}}\label{\detokenize{rozdzial_2/rozdzial_27/rozdzial_1::doc}}
\sphinxAtStartPar
Rozważając bezpieczeństwo systemów bazodanowych, należy spojrzeć na problem warstwowo. Najniższą, a zarazem najbardziej fundamentalną warstwą jest ochrona samej infrastruktury sprzętowej i sieciowej, na której operuje silnik bazy danych (DBMS). Zanim przejdziemy do szczegółowych konfiguracji silnika, konieczne jest zdefiniowanie nadrzędnych celów bezpieczeństwa oraz zapewnienie izolacji środowiska uruchomieniowego.


\subsubsection{Triada CIA w środowisku bazodanowym}
\label{\detokenize{rozdzial_2/rozdzial_27/rozdzial_1:triada-cia-w-srodowisku-bazodanowym}}
\sphinxAtStartPar
Model \sphinxstylestrong{CIA} (ang. \sphinxstyleemphasis{Confidentiality, Integrity, Availability}) to klasyczny paradygmat bezpieczeństwa informacji. W kontekście relacyjnych i nierelacyjnych baz danych każdy z tych elementów realizowany jest przez specyficzne mechanizmy:
\begin{itemize}
\item {} 
\sphinxAtStartPar
\sphinxstylestrong{Poufność (Confidentiality):} Gwarantuje, że dane są dostępne wyłącznie dla autoryzowanych podmiotów. W bazach danych osiąga się to poprzez ścisłą kontrolę dostępu (np. RBAC), szyfrowanie danych w spoczynku (TDE \sphinxhyphen{} \sphinxstyleemphasis{Transparent Data Encryption}), maskowanie danych wrażliwych oraz szyfrowanie połączeń sieciowych (TLS).

\item {} 
\sphinxAtStartPar
\sphinxstylestrong{Integralność (Integrity):} Zapewnia spójność i poprawność danych, chroniąc je przed nieautoryzowaną modyfikacją. Silniki bazodanowe realizują ten postulat natywnie poprzez właściwości ACID (szczególnie \sphinxstyleemphasis{Consistency}), więzy integralności (klucze obce, ograniczenia \sphinxstyleemphasis{CHECK}), a z punktu widzenia bezpieczeństwa \textendash{} poprzez audytowanie DML (Data Manipulation Language) i mechanizmy wykrywania anomalii.

\item {} 
\sphinxAtStartPar
\sphinxstylestrong{Dostępność (Availability):} Oznacza pewność, że system odpowie na żądanie w akceptowalnym czasie. Zagrożeniem dla dostępności są awarie sprzętowe oraz ataki DoS. Obroną na poziomie bazy danych są klastry wysokiej dostępności (HA), replikacja (np. \sphinxstyleemphasis{Master\sphinxhyphen{}Slave}, \sphinxstyleemphasis{Multi\sphinxhyphen{}Master}), partycjonowanie oraz odpowiednio zaprojektowane mechanizmy Disaster Recovery.

\end{itemize}


\subsubsection{Izolacja sieciowa i architektura chmurowa (VPC)}
\label{\detokenize{rozdzial_2/rozdzial_27/rozdzial_1:izolacja-sieciowa-i-architektura-chmurowa-vpc}}
\sphinxAtStartPar
Nawet najlepiej skonfigurowany system uprawnień bazodanowych traci na znaczeniu, jeśli instancja wystawiona jest bezpośrednio do publicznej sieci Internet. Standardem inżynieryjnym w nowoczesnych wdrożeniach (zarówno on\sphinxhyphen{}premise, jak i cloud) jest głęboka separacja sieciowa.

\sphinxAtStartPar
W środowiskach chmurowych (np. AWS, Azure, GCP) bazy danych umieszcza się w dedykowanych chmurach prywatnych (\sphinxstylestrong{VPC} \sphinxhyphen{} \sphinxstyleemphasis{Virtual Private Cloud}). Dobrą praktyką jest wdrożenie architektury wielowarstwowej:
\begin{enumerate}
\sphinxsetlistlabels{\arabic}{enumi}{enumii}{}{.}%
\item {} 
\sphinxAtStartPar
\sphinxstylestrong{Warstwa publiczna (Public Subnet):} Zawiera load balancery i serwery webowe (np. Nginx), do których dostęp ma użytkownik końcowy.

\item {} 
\sphinxAtStartPar
\sphinxstylestrong{Warstwa aplikacji (Private Subnet):} Środowisko uruchomieniowe backendu (np. Node.js, Spring Boot).

\item {} 
\sphinxAtStartPar
\sphinxstylestrong{Warstwa danych (Database Subnet):} Najgłębiej ukryta podsieć prywatna bez routingu do Internetu (brak \sphinxstyleemphasis{Internet Gateway}). Dostęp do niej ma wyłącznie warstwa aplikacji poprzez ściśle określone reguły \sphinxstyleemphasis{Security Groups} lub \sphinxstyleemphasis{Network ACLs} (dopuszczające ruch tylko na konkretnym porcie, np. 5432 dla PostgreSQL, i tylko ze znanych adresów IP warstwy aplikacyjnej).

\end{enumerate}

\sphinxAtStartPar
Takie podejście drastycznie zmniejsza powierzchnię ataku (ang. \sphinxstyleemphasis{attack surface}), eliminując ryzyko bezpośrednich ataków typu \sphinxstyleemphasis{brute\sphinxhyphen{}force} na porty bazy danych z zewnątrz.


\subsubsection{Konteneryzacja a bezpieczeństwo bazy}
\label{\detokenize{rozdzial_2/rozdzial_27/rozdzial_1:konteneryzacja-a-bezpieczenstwo-bazy}}
\sphinxAtStartPar
Uruchamianie baz danych w kontenerach (np. Docker, Kubernetes) stało się powszechne, jednak niesie ze sobą specyficzne wektory zagrożeń. Kontenery dzielą jądro (kernel) z systemem hosta, co oznacza, że kompromitacja bazy danych może prowadzić do ucieczki z kontenera (ang. \sphinxstyleemphasis{container breakout}).

\sphinxAtStartPar
Aby zminimalizować to ryzyko, stosuje się następujące praktyki:
\begin{itemize}
\item {} 
\sphinxAtStartPar
\sphinxstylestrong{Brak uprawnień roota:} Procesy DBMS (np. \sphinxcode{\sphinxupquote{postgres}} lub \sphinxcode{\sphinxupquote{mysqld}}) nigdy nie powinny działać jako użytkownik \sphinxcode{\sphinxupquote{root}} wewnątrz kontenera. Standardowe obrazy często wymagają jawnego zdefiniowania dyrektywy \sphinxcode{\sphinxupquote{USER}} w pliku Dockerfile.

\item {} 
\sphinxAtStartPar
\sphinxstylestrong{Ograniczenia zasobów (Cgroups):} Podatność silnika bazy na ataki typu \sphinxstyleemphasis{Denial of Service} można złagodzić na poziomie infrastruktury, nakładając twarde limity na zużycie CPU i pamięci RAM przez dany kontener. Zapobiega to sytuacji, w której przeciążona baza kładzie cały serwer hosta (Resource Exhaustion).

\item {} 
\sphinxAtStartPar
\sphinxstylestrong{Read\sphinxhyphen{}Only Root Filesystem:} Kontener bazy danych powinien mieć zamontowany system plików w trybie tylko do odczytu, z wyjątkiem ściśle zdefiniowanych wolumenów (ang. \sphinxstyleemphasis{volumes}) przeznaczonych na pliki danych (np. \sphinxcode{\sphinxupquote{/var/lib/postgresql/data}}). Uniemożliwia to atakującemu wstrzyknięcie złośliwych skryptów do katalogów systemowych po udanym wykorzystaniu luki aplikacyjnej.

\end{itemize}

\sphinxstepscope


\subsection{2. Ataki na bazy danych i luki w zabezpieczeniach}
\label{\detokenize{rozdzial_2/rozdzial_27/rozdzial_2:ataki-na-bazy-danych-i-luki-w-zabezpieczeniach}}\label{\detokenize{rozdzial_2/rozdzial_27/rozdzial_2::doc}}
\sphinxAtStartPar
Ponieważ bazy danych przechowują najbardziej krytyczne informacje organizacji, stanowią główny cel cyberataków. Według zestawień takich jak OWASP Top 10, podatności związane ze wstrzykiwaniem kodu (Injection) oraz błędną konfiguracją od lat utrzymują się w czołówce zagrożeń. Zrozumienie mechaniki tych ataków jest kluczowe do zaprojektowania skutecznych mechanizmów obronnych w warstwie aplikacyjnej i bazodanowej.


\subsubsection{Wstrzykiwanie kodu SQL (SQL Injection)}
\label{\detokenize{rozdzial_2/rozdzial_27/rozdzial_2:wstrzykiwanie-kodu-sql-sql-injection}}
\sphinxAtStartPar
SQL Injection (SQLi) to podatność polegająca na możliwości modyfikacji zapytania SQL poprzez nieodpowiednio zwalidowane dane wejściowe. Pozwala to atakującemu na wykonanie nieautoryzowanych operacji na bazie danych.

\sphinxAtStartPar
Klasyczny przykład podatności (język PHP):

\begin{sphinxVerbatim}[commandchars=\\\{\}]
\PYG{x}{\PYGZdl{}username = \PYGZdl{}\PYGZus{}POST[\PYGZsq{}username\PYGZsq{}];}
\PYG{x}{\PYGZdl{}query = \PYGZdq{}SELECT * FROM users WHERE username = \PYGZsq{}\PYGZdl{}username\PYGZsq{}\PYGZdq{};}
\PYG{x}{\PYGZdl{}db\PYGZhy{}\PYGZgt{}execute(\PYGZdl{}query);}
\end{sphinxVerbatim}

\sphinxAtStartPar
Jeśli atakujący jako parametr \sphinxcode{\sphinxupquote{username}} przekaże ciąg \sphinxcode{\sphinxupquote{\textquotesingle{} OR \textquotesingle{}1\textquotesingle{}=\textquotesingle{}1}}, zapytanie przyjmie postać:

\begin{sphinxVerbatim}[commandchars=\\\{\}]
\PYG{k}{SELECT}\PYG{+w}{ }\PYG{o}{*}\PYG{+w}{ }\PYG{k}{FROM}\PYG{+w}{ }\PYG{n}{users}\PYG{+w}{ }\PYG{k}{WHERE}\PYG{+w}{ }\PYG{n}{username}\PYG{+w}{ }\PYG{o}{=}\PYG{+w}{ }\PYG{l+s+s1}{\PYGZsq{}\PYGZsq{}}\PYG{+w}{ }\PYG{k}{OR}\PYG{+w}{ }\PYG{l+s+s1}{\PYGZsq{}1\PYGZsq{}}\PYG{o}{=}\PYG{l+s+s1}{\PYGZsq{}1\PYGZsq{}}\PYG{p}{;}
\end{sphinxVerbatim}

\sphinxAtStartPar
Ponieważ warunek \sphinxcode{\sphinxupquote{\textquotesingle{}1\textquotesingle{}=\textquotesingle{}1\textquotesingle{}}} jest zawsze prawdziwy, silnik bazy danych zwróci wszystkie rekordy z tabeli \sphinxcode{\sphinxupquote{users}}, omijając mechanizm uwierzytelniania.

\sphinxAtStartPar
\sphinxstylestrong{Rodzaje ataków SQLi:}
\begin{enumerate}
\sphinxsetlistlabels{\arabic}{enumi}{enumii}{}{.}%
\item {} 
\sphinxAtStartPar
\sphinxstylestrong{In\sphinxhyphen{}Band SQLi (Classic):} Atakujący używa tego samego kanału komunikacyjnego do przeprowadzenia ataku i zebrania wyników. Najpopularniejsze to \sphinxstyleemphasis{Error\sphinxhyphen{}based SQLi} (wymuszanie błędów silnika w celu wycieku metadanych, np. struktury tabel) oraz \sphinxstyleemphasis{Union\sphinxhyphen{}based SQLi} (użycie operatora \sphinxcode{\sphinxupquote{UNION}} do dołączenia złośliwych zapytań do oryginalnego).

\item {} 
\sphinxAtStartPar
\sphinxstylestrong{Inferential (Blind) SQLi:} Występuje, gdy aplikacja nie zwraca komunikatów o błędach ani bezpośrednich wyników zapytań. Atakujący musi wnioskować o strukturze danych na podstawie zachowania aplikacji.
* \sphinxstylestrong{Boolean\sphinxhyphen{}based:} Wysyłanie zapytań warunkowych i obserwowanie, czy strona ładuje się inaczej.
* \sphinxstylestrong{Time\sphinxhyphen{}based:} Zmuszanie bazy do wstrzymania działania (np. funkcjami \sphinxcode{\sphinxupquote{pg\_sleep()}} w PostgreSQL lub \sphinxcode{\sphinxupquote{WAITFOR DELAY}} w SQL Server), aby potwierdzić obecność podatności.

\item {} 
\sphinxAtStartPar
\sphinxstylestrong{Out\sphinxhyphen{}of\sphinxhyphen{}Band SQLi:} Wykorzystywane rzadziej, opiera się na wymuszeniu przez silnik bazy danych żądania HTTP lub DNS do zewnętrznego serwera kontrolowanego przez atakującego.

\end{enumerate}

\sphinxAtStartPar
\sphinxstylestrong{Mitygacja:}
Podstawową i najskuteczniejszą metodą obrony przed SQLi jest korzystanie z \sphinxstylestrong{zapytań parametryzowanych (Prepared Statements)} lub nowoczesnych systemów ORM (Object\sphinxhyphen{}Relational Mapping), które oddzielają składnię zapytań od danych wejściowych.


\subsubsection{Ataki NoSQL Injection}
\label{\detokenize{rozdzial_2/rozdzial_27/rozdzial_2:ataki-nosql-injection}}
\sphinxAtStartPar
Wraz ze wzrostem popularności nierelacyjnych baz danych (np. MongoDB, CouchDB), pojawił się mit o ich całkowitej odporności na wstrzykiwanie kodu. Choć nie używają one tradycyjnego dialektu SQL, aplikacje z nich korzystające nadal mogą być podatne na modyfikację logiki zapytań.

\sphinxAtStartPar
Ataki te często bazują na wstrzykiwaniu operatorów bazy danych do zapytań, które aplikacja interpretuje jako obiekty JSON.

\sphinxAtStartPar
\sphinxstylestrong{Przykład dla MongoDB:}
Aplikacja Node.js oczekuje poświadczeń przekazanych w formacie JSON:

\begin{sphinxVerbatim}[commandchars=\\\{\}]
\PYG{n+nx}{db}\PYG{p}{.}\PYG{n+nx}{collection}\PYG{p}{(}\PYG{l+s+s1}{\PYGZsq{}users\PYGZsq{}}\PYG{p}{)}\PYG{p}{.}\PYG{n+nx}{find}\PYG{p}{(}\PYG{p}{\PYGZob{}}
\PYG{+w}{    }\PYG{n+nx}{username}\PYG{o}{:}\PYG{+w}{ }\PYG{n+nx}{req}\PYG{p}{.}\PYG{n+nx}{body}\PYG{p}{.}\PYG{n+nx}{username}\PYG{p}{,}
\PYG{+w}{    }\PYG{n+nx}{password}\PYG{o}{:}\PYG{+w}{ }\PYG{n+nx}{req}\PYG{p}{.}\PYG{n+nx}{body}\PYG{p}{.}\PYG{n+nx}{password}
\PYG{p}{\PYGZcb{}}\PYG{p}{)}\PYG{p}{;}
\end{sphinxVerbatim}

\sphinxAtStartPar
Atakujący wysyła zmodyfikowany payload z użyciem operatora logicznego \sphinxcode{\sphinxupquote{\$ne}} (Not Equal):

\begin{sphinxVerbatim}[commandchars=\\\{\}]
\PYG{p}{\PYGZob{}}
\PYG{+w}{    }\PYG{n+nt}{\PYGZdq{}username\PYGZdq{}}\PYG{p}{:}\PYG{+w}{ }\PYG{p}{\PYGZob{}}\PYG{n+nt}{\PYGZdq{}\PYGZdl{}ne\PYGZdq{}}\PYG{p}{:}\PYG{+w}{ }\PYG{k+kc}{null}\PYG{p}{\PYGZcb{},}
\PYG{+w}{    }\PYG{n+nt}{\PYGZdq{}password\PYGZdq{}}\PYG{p}{:}\PYG{+w}{ }\PYG{p}{\PYGZob{}}\PYG{n+nt}{\PYGZdq{}\PYGZdl{}ne\PYGZdq{}}\PYG{p}{:}\PYG{+w}{ }\PYG{k+kc}{null}\PYG{p}{\PYGZcb{}}
\PYG{p}{\PYGZcb{}}
\end{sphinxVerbatim}

\sphinxAtStartPar
W rezultacie silnik MongoDB wykonuje zapytanie szukające użytkownika, którego login i hasło \sphinxstyleemphasis{nie są nullem}. Zapytanie zwraca pierwszy pasujący rekord (często administratora), umożliwiając logowanie bez znajomości hasła.


\subsubsection{Przepełnienie bufora (Buffer Overflow) w DBMS}
\label{\detokenize{rozdzial_2/rozdzial_27/rozdzial_2:przepelnienie-bufora-buffer-overflow-w-dbms}}
\sphinxAtStartPar
Silniki baz danych (np. MySQL, PostgreSQL, Oracle) to skomplikowane aplikacje napisane najczęściej w językach C lub C++, co oznacza, że są potencjalnie narażone na błędy zarządzania pamięcią.

\sphinxAtStartPar
\sphinxstyleemphasis{Buffer Overflow} w kontekście baz danych występuje, gdy atakujący przesyła nieproporcjonalnie duży ciąg danych do bufora o stałej wielkości (np. w parametrze funkcji wbudowanej, mechanizmie autoryzacyjnym lub procedurze składowanej). Gdy dane wykraczają poza zaalokowany obszar pamięci, mogą nadpisać sąsiadujące obszary, w tym wskaźnik powrotu instrukcji (Instruction Pointer).

\sphinxAtStartPar
Może to prowadzić do:
1. \sphinxstylestrong{Awarji silnika bazy danych (Crash)} \textendash{} przerywając ciągłość działania.
2. \sphinxstylestrong{Wykonania dowolnego kodu (RCE \sphinxhyphen{} Remote Code Execution)} \textendash{} uruchomienia shellcode’u na serwerze z uprawnieniami procesu bazy danych, co stanowi bezpośrednie wejście do kompromitacji hosta (powiązane z ryzykiem opisanym w rozdziale o konteneryzacji).


\subsubsection{Odmowa Usługi (DoS) na poziomie bazy danych}
\label{\detokenize{rozdzial_2/rozdzial_27/rozdzial_2:odmowa-uslugi-dos-na-poziomie-bazy-danych}}
\sphinxAtStartPar
Ataki Denial of Service na bazy danych różnią się od klasycznych ataków sieciowych. Mają one charakter aplikacyjny i celują w wyczerpanie zasobów sprzętowych serwera (CPU, RAM, I/O) lub limitów samego oprogramowania.
\begin{itemize}
\item {} 
\sphinxAtStartPar
\sphinxstylestrong{Wyczerpanie puli połączeń (Connection Exhaustion):} Każdy DBMS ma zdefiniowany limit maksymalnej liczby jednoczesnych połączeń (np. parametr \sphinxcode{\sphinxupquote{max\_connections}} w PostgreSQL). Atakujący, często wykorzystując luki w warstwie aplikacji, może zainicjować tysiące zawieszonych lub bardzo powolnych sesji bazy, uniemożliwiając logowanie legalnym użytkownikom.

\item {} 
\sphinxAtStartPar
\sphinxstylestrong{Asymetryczne obciążenie (Query\sphinxhyphen{}based DoS):} Atakujący wstrzykuje i uruchamia zapytania, które są syntaktycznie poprawne, ale drastycznie nieoptymalne kosztowo. Przykładem może być wymuszenie iloczynu kartezjańskiego (CROSS JOIN) na bardzo dużych tabelach bez odpowiednich indeksów, co może zablokować procesor serwera na kilka minut i doprowadzić do blokady całej aplikacji (Deadlock/Timeouts).

\end{itemize}


\subsubsection{Automatyzacja ataków: Narzędzia audytowe}
\label{\detokenize{rozdzial_2/rozdzial_27/rozdzial_2:automatyzacja-atakow-narzedzia-audytowe}}
\sphinxAtStartPar
W procesach testów penetracyjnych baz danych, analitycy bezpieczeństwa korzystają ze zautomatyzowanych frameworków. Najpopularniejszym z nich jest \sphinxstylestrong{sqlmap} \textendash{} potężne narzędzie open\sphinxhyphen{}source automatyzujące proces wykrywania podatności SQLi i przejmowania kontroli nad serwerami bazodanowymi.

\sphinxAtStartPar
W typowym scenariuszu audytu, uruchomienie polecenia:

\begin{sphinxVerbatim}[commandchars=\\\{\}]
sqlmap\PYG{+w}{ }\PYGZhy{}u\PYG{+w}{ }\PYG{l+s+s2}{\PYGZdq{}http://podatna\PYGZhy{}aplikacja.local/item.php?id=1\PYGZdq{}}\PYG{+w}{ }\PYGZhy{}\PYGZhy{}dbs
\end{sphinxVerbatim}

\sphinxAtStartPar
pozwoli narzędziu na wykrycie typu wstrzyknięcia (np. czasowe lub błędowe), zidentyfikowanie używanego silnika DBMS oraz \textendash{} w razie sukcesu \textendash{} wyliczenie wszystkich baz danych dostępnych na serwerze (enumeracja instancji).

\sphinxstepscope


\chapter{Rozdział 3: Projektowanie Bazy Danych: System Zarządzania Biblioteką}
\label{\detokenize{rozdzial_3/index:rozdzial-3-projektowanie-bazy-danych-system-zarzadzania-biblioteka}}\label{\detokenize{rozdzial_3/index::doc}}\begin{quote}\begin{description}
\sphinxlineitem{Autorzy}\begin{enumerate}
\sphinxsetlistlabels{\arabic}{enumi}{enumii}{}{.}%
\item {} 
\sphinxAtStartPar
Paweł Łoćwin

\item {} 
\sphinxAtStartPar
Paweł Łosowski

\end{enumerate}

\end{description}\end{quote}


\section{Wybór zagadnienia, opis procesów i danych}
\label{\detokenize{rozdzial_3/index:wybor-zagadnienia-opis-procesow-i-danych}}

\subsection{Wybrane zagadnienie:}
\label{\detokenize{rozdzial_3/index:wybrane-zagadnienie}}
\sphinxAtStartPar
System zarządzania wypożyczeniami w bibliotece.
Projekt obejmuje obsługę bazy czytelników, katalogu zbiorów bibliotecznych (z uwzględnieniem autorów i kategorii literackich), a także pełen proces ewidencji wypożyczeń, zwrotów oraz historii operacji każdego użytkownika.


\subsection{Opis procesów i więzy integralności:}
\label{\detokenize{rozdzial_3/index:opis-procesow-i-wiezy-integralnosci}}
\sphinxAtStartPar
Głównym procesem w systemie jest obieg książki między biblioteką a czytelnikiem.
\begin{itemize}
\item {} 
\sphinxAtStartPar
\sphinxstylestrong{Rejestracja:} Czytelnik zapisuje się do biblioteki, podając swoje dane osobowe oraz adresowe. System automatycznie rejestruje datę zapisu.

\item {} 
\sphinxAtStartPar
\sphinxstylestrong{Katalogowanie:} Książki są kategoryzowane (np. Fantastyka, Kryminał, Literatura faktu) i przypisane do konkretnych autorów. Każdy fizyczny egzemplarz książki posiada swój unikalny identyfikator, co umożliwia śledzenie niezależnie każdego egzemplarza.

\item {} 
\sphinxAtStartPar
\sphinxstylestrong{Wypożyczenia i Zwroty:} Proces wypożyczenia łączy czytelnika z konkretnym egzemplarzem książki w czasie. Rejestrowana jest data wypożyczenia. System zakłada konieczność ewidencji dokładnej daty zwrotu rzeczywistego, co pozwala na analizy przeterminowanych wypożyczeń i wzorców korzystania z kolekcji.

\item {} 
\sphinxAtStartPar
\sphinxstylestrong{Statystyki (Więzy integralności):} Pola takie jak “aktualne wypożyczenia” i “suma wypożyczeń” z pierwotnego prototypu są wartościami wyliczalnymi (pochodnymi) na podstawie historii transakcji i są usuwane ze schematu normalnego.

\end{itemize}


\subsection{Wykaz gromadzonych danych:}
\label{\detokenize{rozdzial_3/index:wykaz-gromadzonych-danych}}\begin{itemize}
\item {} 
\sphinxAtStartPar
\sphinxstylestrong{Dane czytelnika:} Imię, Nazwisko, Ulica, Kod pocztowy, Miasto, Data zapisu.

\item {} 
\sphinxAtStartPar
\sphinxstylestrong{Dane książki:} Tytuł, Rok wydania, Dostępność.

\item {} 
\sphinxAtStartPar
\sphinxstylestrong{Dane autora:} Imię, Nazwisko.

\item {} 
\sphinxAtStartPar
\sphinxstylestrong{Dane kategorii:} Nazwa gatunku literackiego.

\item {} 
\sphinxAtStartPar
\sphinxstylestrong{Dane transakcyjne:} Data wypożyczenia, Data zwrotu (rzeczywista).

\end{itemize}


\section{Prototyp CSV}
\label{\detokenize{rozdzial_3/index:prototyp-csv}}
\sphinxAtStartPar
Aby zweryfikować kompletność przetwarzanych informacji, przygotowano “płaską” (nieznormalizowaną) reprezentację danych dla operacji wypożyczenia, bazującą na pierwotnych wytycznych.

\begin{sphinxVerbatim}[commandchars=\\\{\}]
\PYG{n}{Imie}\PYG{p}{,}\PYG{n}{Nazwisko}\PYG{p}{,}\PYG{n}{Ulica}\PYG{p}{,}\PYG{n}{Kod\PYGZus{}Pocztowy}\PYG{p}{,}\PYG{n}{Miasto}\PYG{p}{,}\PYG{n}{Data\PYGZus{}Zapisu}\PYG{p}{,}\PYG{n}{Tytul\PYGZus{}Ksiazki}\PYG{p}{,}\PYG{n}{Imie\PYGZus{}Autora}\PYG{p}{,}\PYG{n}{Nazwisko\PYGZus{}Autora}\PYG{p}{,}\PYG{n}{Kategoria}\PYG{p}{,}\PYG{n}{Data\PYGZus{}Wypozyczenia}\PYG{p}{,}\PYG{n}{Data\PYGZus{}Zwrotu}
\PYG{n}{Jan}\PYG{p}{,}\PYG{n}{Kowalski}\PYG{p}{,}\PYG{n}{Kwiatowa} \PYG{l+m+mi}{1}\PYG{p}{,}\PYG{l+m+mi}{00}\PYG{o}{\PYGZhy{}}\PYG{l+m+mi}{001}\PYG{p}{,}\PYG{n}{Warszawa}\PYG{p}{,}\PYG{l+m+mi}{2024}\PYG{o}{\PYGZhy{}}\PYG{l+m+mi}{01}\PYG{o}{\PYGZhy{}}\PYG{l+m+mi}{15}\PYG{p}{,}\PYG{n}{Wiedźmin}\PYG{p}{,}\PYG{n}{Andrzej}\PYG{p}{,}\PYG{n}{Sapkowski}\PYG{p}{,}\PYG{n}{Fantastyka}\PYG{p}{,}\PYG{l+m+mi}{2024}\PYG{o}{\PYGZhy{}}\PYG{l+m+mi}{03}\PYG{o}{\PYGZhy{}}\PYG{l+m+mi}{01}\PYG{p}{,}\PYG{l+m+mi}{2024}\PYG{o}{\PYGZhy{}}\PYG{l+m+mi}{03}\PYG{o}{\PYGZhy{}}\PYG{l+m+mi}{15}
\PYG{n}{Jan}\PYG{p}{,}\PYG{n}{Kowalski}\PYG{p}{,}\PYG{n}{Kwiatowa} \PYG{l+m+mi}{1}\PYG{p}{,}\PYG{l+m+mi}{00}\PYG{o}{\PYGZhy{}}\PYG{l+m+mi}{001}\PYG{p}{,}\PYG{n}{Warszawa}\PYG{p}{,}\PYG{l+m+mi}{2024}\PYG{o}{\PYGZhy{}}\PYG{l+m+mi}{01}\PYG{o}{\PYGZhy{}}\PYG{l+m+mi}{15}\PYG{p}{,}\PYG{n}{Pan} \PYG{n}{Tadeusz}\PYG{p}{,}\PYG{n}{Adam}\PYG{p}{,}\PYG{n}{Mickiewicz}\PYG{p}{,}\PYG{n}{Poezja}\PYG{p}{,}\PYG{l+m+mi}{2024}\PYG{o}{\PYGZhy{}}\PYG{l+m+mi}{03}\PYG{o}{\PYGZhy{}}\PYG{l+m+mi}{10}\PYG{p}{,}
\PYG{n}{Anna}\PYG{p}{,}\PYG{n}{Nowak}\PYG{p}{,}\PYG{n}{Polna} \PYG{l+m+mi}{2}\PYG{p}{,}\PYG{l+m+mi}{31}\PYG{o}{\PYGZhy{}}\PYG{l+m+mi}{444}\PYG{p}{,}\PYG{n}{Kraków}\PYG{p}{,}\PYG{l+m+mi}{2023}\PYG{o}{\PYGZhy{}}\PYG{l+m+mi}{11}\PYG{o}{\PYGZhy{}}\PYG{l+m+mi}{05}\PYG{p}{,}\PYG{n}{Lśnienie}\PYG{p}{,}\PYG{n}{Stephen}\PYG{p}{,}\PYG{n}{King}\PYG{p}{,}\PYG{n}{Horror}\PYG{p}{,}\PYG{l+m+mi}{2024}\PYG{o}{\PYGZhy{}}\PYG{l+m+mi}{02}\PYG{o}{\PYGZhy{}}\PYG{l+m+mi}{20}\PYG{p}{,}\PYG{l+m+mi}{2024}\PYG{o}{\PYGZhy{}}\PYG{l+m+mi}{03}\PYG{o}{\PYGZhy{}}\PYG{l+m+mi}{05}
\end{sphinxVerbatim}


\section{Model Konceptualny (Pojęciowy)}
\label{\detokenize{rozdzial_3/index:model-konceptualny-pojeciowy}}
\sphinxAtStartPar
Na podstawie analizy procesów i zebranych danych opracowano model pojęciowy, identyfikując obiekty, ich cechy oraz powiązania.


\subsection{Zidentyfikowane encje}
\label{\detokenize{rozdzial_3/index:zidentyfikowane-encje}}\begin{itemize}
\item {} 
\sphinxAtStartPar
\sphinxstylestrong{Czytelnik} \sphinxhyphen{} osoba zapisana do biblioteki.

\item {} 
\sphinxAtStartPar
\sphinxstylestrong{Autor} \sphinxhyphen{} twórca dzieła literackiego.

\item {} 
\sphinxAtStartPar
\sphinxstylestrong{Książka} \sphinxhyphen{} oferowana pozycja z inwentarza biblioteki.

\item {} 
\sphinxAtStartPar
\sphinxstylestrong{Kategoria} \sphinxhyphen{} sekcja grupująca gatunkowo księgozbiór.

\item {} 
\sphinxAtStartPar
\sphinxstylestrong{Wypożyczenie} \sphinxhyphen{} zdarzenie potwierdzające fakt wydania książki czytelnikowi.

\end{itemize}


\subsection{Zdefiniowane atrybuty/własności}
\label{\detokenize{rozdzial_3/index:zdefiniowane-atrybuty-wlasnosci}}\begin{itemize}
\item {} 
\sphinxAtStartPar
\sphinxstylestrong{Czytelnik:} Imię, Nazwisko, Ulica, Kod pocztowy, Miasto, Data zapisu.

\item {} 
\sphinxAtStartPar
\sphinxstylestrong{Autor:} Imię, Nazwisko.

\item {} 
\sphinxAtStartPar
\sphinxstylestrong{Książka:} Tytuł, Rok wydania.

\item {} 
\sphinxAtStartPar
\sphinxstylestrong{Kategoria:} Nazwa kategorii.

\item {} 
\sphinxAtStartPar
\sphinxstylestrong{Wypożyczenie:} Data wypożyczenia, Data zwrotu.

\end{itemize}


\subsection{Opis związków}
\label{\detokenize{rozdzial_3/index:opis-zwiazkow}}\begin{itemize}
\item {} 
\sphinxAtStartPar
\sphinxstylestrong{Autor \textendash{} Książka (1:N):} Jeden autor może napisać wiele książek. (Dla uproszczenia modelu zakładamy głównego autora).

\item {} 
\sphinxAtStartPar
\sphinxstylestrong{Kategoria \textendash{} Książka (1:N):} Kategoria zawiera wiele tytułów.

\item {} 
\sphinxAtStartPar
\sphinxstylestrong{Czytelnik \textendash{} Wypożyczenie (1:N):} Czytelnik może dokonać wielu wypożyczeń w historii.

\item {} 
\sphinxAtStartPar
\sphinxstylestrong{Książka \textendash{} Wypożyczenie (1:N):} Jeden egzemplarz książki może być w swojej historii wypożyczany wielokrotnie różnym czytelnikom.

\end{itemize}


\subsection{Identyfikacja encji słabych}
\label{\detokenize{rozdzial_3/index:identyfikacja-encji-slabych}}
\sphinxAtStartPar
Występuje relacja wiele\sphinxhyphen{}do\sphinxhyphen{}wielu (M:N) między Czytelnikiem a Książką (czytelnik czyta wiele książek, książka jest czytana przez wielu czytelników). Związek ten rozwiązywany jest przez encję asocjacyjną \sphinxstylestrong{Wypożyczenia}, która przechowuje wszystkie historyczne i bieżące transakcje między dwoma stronami.
* \sphinxstylestrong{Uzasadnienie:} Byt ten nie istnieje samodzielnie w oderwaniu od konkretnego Czytelnika i konkretnej Książki.


\subsection{Schemat w notacji Chena}
\label{\detokenize{rozdzial_3/index:schemat-w-notacji-chena}}
\noindent\sphinxincludegraphics[width=0.800\linewidth]{{model_konceptualny_biblioteka}.png}


\section{Model logiczny i proces normalizacji}
\label{\detokenize{rozdzial_3/index:model-logiczny-i-proces-normalizacji}}

\subsection{Przebieg procesu normalizacji}
\label{\detokenize{rozdzial_3/index:przebieg-procesu-normalizacji}}
\sphinxAtStartPar
\sphinxstylestrong{Krok 1: Pierwsza Postać Normalna (1NF)}
Wiersze są unikalne, tworzymy sztuczne klucze główne (ID\_Wypozyczenia). Wartości w komórkach są atomowe. Pojawia się duża redundancja danych adresowych i książkowych.

\sphinxAtStartPar
\sphinxstylestrong{Krok 2: Druga Postać Normalna (2NF)}
Rozbijamy płaską strukturę względem częściowych zależności. Oddzielamy encje twarde od operacji.
Tworzymy bazowe tabele: \sphinxcode{\sphinxupquote{Czytelnicy}} oraz \sphinxcode{\sphinxupquote{Ksiazki}}. Powstaje tabela \sphinxcode{\sphinxupquote{Wypozyczenia}} przechowująca klucze obce ID\_Czytelnika oraz ID\_Ksiazki, a także daty transakcji.
\sphinxstyleemphasis{Uwaga:} Dynamiczne atrybuty takie jak “suma wypożyczeń” z pierwotnego zadania są usuwane ze schematu, ponieważ łamią zasady normalizacji \textendash{} można je obliczyć zapytaniem SQL typu \sphinxcode{\sphinxupquote{COUNT(*)}} na zdarzeniach w tabeli \sphinxcode{\sphinxupquote{Wypozyczenia}}.

\sphinxAtStartPar
\sphinxstylestrong{Krok 3: Trzecia Postać Normalna (3NF)}
Eliminacja zależności przechodnich. Z tabeli \sphinxcode{\sphinxupquote{Ksiazki}} wydzielamy powtarzające się nazwy gatunków do tabeli \sphinxcode{\sphinxupquote{Kategorie}} (klucz: ID\_Kategorii) oraz dane twórców do tabeli \sphinxcode{\sphinxupquote{Autorzy}} (klucz: ID\_Autora). Każda tabela reprezentuje teraz niezależną encję o jasno określonym celu biznesowym.


\subsection{Ostateczna struktura tabel (3NF)}
\label{\detokenize{rozdzial_3/index:ostateczna-struktura-tabel-3nf}}\begin{itemize}
\item {} 
\sphinxAtStartPar
\sphinxstylestrong{Czytelnicy:} ID\_Czytelnika (PK), Imie, Nazwisko, Ulica, Kod\_Pocztowy, Miasto, Data\_Zapisu

\item {} 
\sphinxAtStartPar
\sphinxstylestrong{Autorzy:} ID\_Autora (PK), Imie, Nazwisko

\item {} 
\sphinxAtStartPar
\sphinxstylestrong{Kategorie:} ID\_Kategorii (PK), Nazwa\_Kategorii

\item {} 
\sphinxAtStartPar
\sphinxstylestrong{Ksiazki:} ID\_Ksiazki (PK), ID\_Autora (FK), ID\_Kategorii (FK), Tytul, Rok\_Wydania

\item {} 
\sphinxAtStartPar
\sphinxstylestrong{Wypozyczenia:} ID\_Wypozyczenia (PK), ID\_Czytelnika (FK), ID\_Ksiazki (FK), Data\_Wypozyczenia, Data\_Zwrotu

\end{itemize}


\subsection{Diagram ERD (Model Logiczny)}
\label{\detokenize{rozdzial_3/index:diagram-erd-model-logiczny}}
\noindent\sphinxincludegraphics[width=0.800\linewidth]{{erd_logiczny_biblioteka}.png}


\section{Model fizyczny bazy danych}
\label{\detokenize{rozdzial_3/index:model-fizyczny-bazy-danych}}
\sphinxAtStartPar
Różnice implementacyjne między modelami fizycznymi wynikają wprost z silników RDBMS \textendash{} ograniczeń typowania SQLite oraz bogatych możliwości deklaratywnych PostgreSQL.


\subsection{Model fizyczny dla środowiska SQLite}
\label{\detokenize{rozdzial_3/index:model-fizyczny-dla-srodowiska-sqlite}}
\sphinxAtStartPar
Z uwagi na okrojony zestaw typów, daty są mapowane jako TEXT.
\begin{itemize}
\item {} 
\sphinxAtStartPar
\sphinxstylestrong{Czytelnicy:} ID\_Czytelnika : INTEGER PRIMARY KEY, Imie : TEXT, Nazwisko : TEXT, Ulica : TEXT, Kod\_Pocztowy : TEXT, Miasto : TEXT, Data\_Zapisu : TEXT

\item {} 
\sphinxAtStartPar
\sphinxstylestrong{Autorzy:} ID\_Autora : INTEGER PRIMARY KEY, Imie : TEXT, Nazwisko : TEXT

\item {} 
\sphinxAtStartPar
\sphinxstylestrong{Kategorie:} ID\_Kategorii : INTEGER PRIMARY KEY, Nazwa\_Kategorii : TEXT

\item {} 
\sphinxAtStartPar
\sphinxstylestrong{Ksiazki:} ID\_Ksiazki : INTEGER PRIMARY KEY, ID\_Autora : INTEGER, ID\_Kategorii : INTEGER, Tytul : TEXT, Rok\_Wydania : INTEGER

\item {} 
\sphinxAtStartPar
\sphinxstylestrong{Wypozyczenia:} ID\_Wypozyczenia : INTEGER PRIMARY KEY, ID\_Czytelnika : INTEGER, ID\_Ksiazki : INTEGER, Data\_Wypozyczenia : TEXT, Data\_Zwrotu : TEXT

\end{itemize}

\noindent\sphinxincludegraphics[width=0.800\linewidth]{{fizyczny_sqlite_biblioteka}.png}


\subsection{Model fizyczny dla środowiska PostgreSQL}
\label{\detokenize{rozdzial_3/index:model-fizyczny-dla-srodowiska-postgresql}}
\sphinxAtStartPar
PostgreSQL umożliwia zastosowanie precyzyjnych i natywnych typów, w tym rygorystycznych typów daty (DATE) oraz optymalizacji pamięciowej dla ciągów znaków (VARCHAR).
\begin{itemize}
\item {} 
\sphinxAtStartPar
\sphinxstylestrong{Czytelnicy:} ID\_Czytelnika : SERIAL PRIMARY KEY, Imie : VARCHAR(50), Nazwisko : VARCHAR(50), Ulica : VARCHAR(100), Kod\_Pocztowy : VARCHAR(6), Miasto : VARCHAR(50), Data\_Zapisu : DATE DEFAULT CURRENT\_DATE

\item {} 
\sphinxAtStartPar
\sphinxstylestrong{Autorzy:} ID\_Autora : SERIAL PRIMARY KEY, Imie : VARCHAR(50), Nazwisko : VARCHAR(50)

\item {} 
\sphinxAtStartPar
\sphinxstylestrong{Kategorie:} ID\_Kategorii : SERIAL PRIMARY KEY, Nazwa\_Kategorii : VARCHAR(50)

\item {} 
\sphinxAtStartPar
\sphinxstylestrong{Ksiazki:} ID\_Ksiazki : SERIAL PRIMARY KEY, ID\_Autora : INTEGER REFERENCES Autorzy, ID\_Kategorii : INTEGER REFERENCES Kategorie, Tytul : VARCHAR(150), Rok\_Wydania : SMALLINT

\item {} 
\sphinxAtStartPar
\sphinxstylestrong{Wypozyczenia:} ID\_Wypozyczenia : SERIAL PRIMARY KEY, ID\_Czytelnika : INTEGER REFERENCES Czytelnicy, ID\_Ksiazki : INTEGER REFERENCES Ksiazki, Data\_Wypozyczenia : DATE DEFAULT CURRENT\_DATE, Data\_Zwrotu : DATE

\end{itemize}

\noindent\sphinxincludegraphics[width=0.800\linewidth]{{fizyczny_postgres_biblioteka.drawio}.png}

\sphinxstepscope


\chapter{Rozdział 4: Implementacja fizyczna i zasilanie bazy danych}
\label{\detokenize{rozdzial_4/index:rozdzial-4-implementacja-fizyczna-i-zasilanie-bazy-danych}}\label{\detokenize{rozdzial_4/index::doc}}\begin{quote}\begin{description}
\sphinxlineitem{Autorzy}\begin{enumerate}
\sphinxsetlistlabels{\arabic}{enumi}{enumii}{}{.}%
\item {} 
\sphinxAtStartPar
Paweł Łoćwin

\item {} 
\sphinxAtStartPar
Paweł Łosowski

\end{enumerate}

\end{description}\end{quote}


\section{Definicja fizyczna bazy danych (Zadanie 1)}
\label{\detokenize{rozdzial_4/index:definicja-fizyczna-bazy-danych-zadanie-1}}
\sphinxAtStartPar
Na podstawie opracowanych wcześniej modeli logicznych przygotowano skrypty DDL (Data Definition Language) dla dwóch docelowych silników relacyjnych baz danych: PostgreSQL oraz SQLite. Kody te przygotowały solidną fundamentę dla skutecznego przechowywania i zarządzania danymi biblioteki.


\subsection{Skrypt dla środowiska PostgreSQL}
\label{\detokenize{rozdzial_4/index:skrypt-dla-srodowiska-postgresql}}
\sphinxAtStartPar
Środowisko PostgreSQL pozwala na wykorzystanie zaawansowanych, natywnych typów danych oraz rygorystyczne egzekwowanie więzów integralności. W poniższym skrypcie kluczowym elementem jest użycie więzów klucza obcego (FOREIGN KEY) w celu zapewnienia spójności referencyalnej między tabelami.

\begin{sphinxVerbatim}[commandchars=\\\{\}]
\PYG{c+c1}{\PYGZhy{}\PYGZhy{} Fragment kodu definicji kluczowych tabel (PostgreSQL)}
\PYG{k}{CREATE}\PYG{+w}{ }\PYG{k}{TABLE}\PYG{+w}{ }\PYG{n}{Ksiazki}\PYG{+w}{ }\PYG{p}{(}
\PYG{+w}{    }\PYG{n}{ID\PYGZus{}Ksiazki}\PYG{+w}{ }\PYG{n+nb}{SERIAL}\PYG{+w}{ }\PYG{k}{PRIMARY}\PYG{+w}{ }\PYG{k}{KEY}\PYG{p}{,}
\PYG{+w}{    }\PYG{n}{ID\PYGZus{}Autora}\PYG{+w}{ }\PYG{n+nb}{INTEGER}\PYG{+w}{ }\PYG{k}{REFERENCES}\PYG{+w}{ }\PYG{n}{Autorzy}\PYG{p}{(}\PYG{n}{ID\PYGZus{}Autora}\PYG{p}{)}\PYG{p}{,}
\PYG{+w}{    }\PYG{n}{ID\PYGZus{}Kategorii}\PYG{+w}{ }\PYG{n+nb}{INTEGER}\PYG{+w}{ }\PYG{k}{REFERENCES}\PYG{+w}{ }\PYG{n}{Kategorie}\PYG{p}{(}\PYG{n}{ID\PYGZus{}Kategorii}\PYG{p}{)}\PYG{p}{,}
\PYG{+w}{    }\PYG{n}{Tytul}\PYG{+w}{ }\PYG{n+nb}{VARCHAR}\PYG{p}{(}\PYG{l+m+mi}{150}\PYG{p}{)}\PYG{+w}{ }\PYG{k}{NOT}\PYG{+w}{ }\PYG{k}{NULL}\PYG{p}{,}
\PYG{+w}{    }\PYG{n}{Rok\PYGZus{}Wydania}\PYG{+w}{ }\PYG{n+nb}{SMALLINT}
\PYG{p}{)}\PYG{p}{;}

\PYG{k}{CREATE}\PYG{+w}{ }\PYG{k}{TABLE}\PYG{+w}{ }\PYG{n}{Wypozyczenia}\PYG{+w}{ }\PYG{p}{(}
\PYG{+w}{    }\PYG{n}{ID\PYGZus{}Wypozyczenia}\PYG{+w}{ }\PYG{n+nb}{SERIAL}\PYG{+w}{ }\PYG{k}{PRIMARY}\PYG{+w}{ }\PYG{k}{KEY}\PYG{p}{,}
\PYG{+w}{    }\PYG{n}{ID\PYGZus{}Czytelnika}\PYG{+w}{ }\PYG{n+nb}{INTEGER}\PYG{+w}{ }\PYG{k}{REFERENCES}\PYG{+w}{ }\PYG{n}{Czytelnicy}\PYG{p}{(}\PYG{n}{ID\PYGZus{}Czytelnika}\PYG{p}{)}\PYG{p}{,}
\PYG{+w}{    }\PYG{n}{ID\PYGZus{}Ksiazki}\PYG{+w}{ }\PYG{n+nb}{INTEGER}\PYG{+w}{ }\PYG{k}{REFERENCES}\PYG{+w}{ }\PYG{n}{Ksiazki}\PYG{p}{(}\PYG{n}{ID\PYGZus{}Ksiazki}\PYG{p}{)}\PYG{p}{,}
\PYG{+w}{    }\PYG{n}{Data\PYGZus{}Wypozyczenia}\PYG{+w}{ }\PYG{n+nb}{DATE}\PYG{+w}{ }\PYG{k}{DEFAULT}\PYG{+w}{ }\PYG{k}{CURRENT\PYGZus{}DATE}\PYG{p}{,}
\PYG{+w}{    }\PYG{n}{Data\PYGZus{}Zwrotu}\PYG{+w}{ }\PYG{n+nb}{DATE}
\PYG{p}{)}\PYG{p}{;}
\end{sphinxVerbatim}


\subsection{Skrypt dla środowiska SQLite}
\label{\detokenize{rozdzial_4/index:skrypt-dla-srodowiska-sqlite}}
\sphinxAtStartPar
W silniku SQLite struktura ulega pewnemu uproszczeniu z powodu mniejszej rygorystyczności typowania oraz bardzo ograniczonej liczby wbudowanych klas typów danych (np. brak odrębnego typu daty). Pomimo ograniczeń, SQLite pozostaje niezawodnym rozwiązaniem do prototypowania i mniejszych aplikacji.

\begin{sphinxVerbatim}[commandchars=\\\{\}]
\PYG{c+c1}{\PYGZhy{}\PYGZhy{} Fragment kodu definicji kluczowych tabel (SQLite)}
\PYG{k}{CREATE}\PYG{+w}{ }\PYG{k}{TABLE}\PYG{+w}{ }\PYG{n}{Wypozyczenia}\PYG{+w}{ }\PYG{p}{(}
\PYG{+w}{    }\PYG{n}{ID\PYGZus{}Wypozyczenia}\PYG{+w}{ }\PYG{n+nb}{INTEGER}\PYG{+w}{ }\PYG{k}{PRIMARY}\PYG{+w}{ }\PYG{k}{KEY}\PYG{+w}{ }\PYG{n}{AUTOINCREMENT}\PYG{p}{,}
\PYG{+w}{    }\PYG{n}{ID\PYGZus{}Czytelnika}\PYG{+w}{ }\PYG{n+nb}{INTEGER}\PYG{p}{,}
\PYG{+w}{    }\PYG{n}{ID\PYGZus{}Ksiazki}\PYG{+w}{ }\PYG{n+nb}{INTEGER}\PYG{p}{,}
\PYG{+w}{    }\PYG{n}{Data\PYGZus{}Wypozyczenia}\PYG{+w}{ }\PYG{n+nb}{TEXT}\PYG{p}{,}
\PYG{+w}{    }\PYG{n}{Data\PYGZus{}Zwrotu}\PYG{+w}{ }\PYG{n+nb}{TEXT}\PYG{p}{,}
\PYG{+w}{    }\PYG{k}{FOREIGN}\PYG{+w}{ }\PYG{k}{KEY}\PYG{p}{(}\PYG{n}{ID\PYGZus{}Czytelnika}\PYG{p}{)}\PYG{+w}{ }\PYG{k}{REFERENCES}\PYG{+w}{ }\PYG{n}{Czytelnicy}\PYG{p}{(}\PYG{n}{ID\PYGZus{}Czytelnika}\PYG{p}{)}\PYG{p}{,}
\PYG{+w}{    }\PYG{k}{FOREIGN}\PYG{+w}{ }\PYG{k}{KEY}\PYG{p}{(}\PYG{n}{ID\PYGZus{}Ksiazki}\PYG{p}{)}\PYG{+w}{ }\PYG{k}{REFERENCES}\PYG{+w}{ }\PYG{n}{Ksiazki}\PYG{p}{(}\PYG{n}{ID\PYGZus{}Ksiazki}\PYG{p}{)}
\PYG{p}{)}\PYG{p}{;}
\end{sphinxVerbatim}


\section{Mechanizm zasilania bazy danych (Zadanie 2)}
\label{\detokenize{rozdzial_4/index:mechanizm-zasilania-bazy-danych-zadanie-2}}
\sphinxAtStartPar
Posiadając gotową strukturę bazy, należało wyposażyć ją w dane demonstracyjne. Przygotowano zbiór danych testowych, które zostały poddane wsadowemu ładowaniu (batch import).


\subsection{Formatyzacja danych (Pliki CSV)}
\label{\detokenize{rozdzial_4/index:formatyzacja-danych-pliki-csv}}
\sphinxAtStartPar
Dane do importu przygotowano w niezależnym formacie płaskim CSV (Comma\sphinxhyphen{}Separated Values). Takie podejście pozwala na łatwą edycję danych poza systemem bazodanowym. Poniżej zaprezentowano wybrane kategorie książek przygotowane w formacie CSV.

\begin{sphinxVerbatim}[commandchars=\\\{\}]
Nazwa\PYGZus{}Kategorii
Fantastyka
Kryminał
Literatura faktu
Poezja
Horror
\end{sphinxVerbatim}


\subsection{Implementacja mechanizmu importu}
\label{\detokenize{rozdzial_4/index:implementacja-mechanizmu-importu}}
\sphinxAtStartPar
Do zaimportowania powyższych danych wybrano mechanizm wprowadzania wielowartościowego oparty na technice \sphinxstylestrong{INSERT (multi\sphinxhyphen{}value/batch)}. Rozwiązanie zrealizowano z wykorzystaniem języka Python wraz z biblioteką psycopg do komunikacji z bazą PostgreSQL.

\sphinxAtStartPar
Wybór tego mechanizmu podyktowany jest kompromisem między uniwersalnością a wydajnością. Użycie metody \sphinxcode{\sphinxupquote{executemany()}} na obiekcie kursora bazy danych pozwala na szybkie wstawienie wielu rekordów bez konieczności wykonywania zapytania dla każdego wiersza osobno.

\sphinxAtStartPar
Poniżej przedstawiono fragment skryptu aplikacyjnego odpowiedzialnego za odczyt z pliku i zasilenie tabel:

\begin{sphinxVerbatim}[commandchars=\\\{\}]
\PYG{k+kn}{import} \PYG{n+nn}{csv}
\PYG{k+kn}{import} \PYG{n+nn}{psycopg}

\PYG{c+c1}{\PYGZsh{} 1. Wczytywanie danych z pliku CSV do struktury iterowalnej (listy krotek)}
\PYG{k}{def} \PYG{n+nf}{wczytaj\PYGZus{}dane\PYGZus{}csv}\PYG{p}{(}\PYG{n}{sciezka\PYGZus{}pliku}\PYG{p}{)}\PYG{p}{:}
    \PYG{n}{dane} \PYG{o}{=} \PYG{p}{[}\PYG{p}{]}
    \PYG{k}{with} \PYG{n+nb}{open}\PYG{p}{(}\PYG{n}{sciezka\PYGZus{}pliku}\PYG{p}{,} \PYG{n}{mode}\PYG{o}{=}\PYG{l+s+s1}{\PYGZsq{}}\PYG{l+s+s1}{r}\PYG{l+s+s1}{\PYGZsq{}}\PYG{p}{,} \PYG{n}{encoding}\PYG{o}{=}\PYG{l+s+s1}{\PYGZsq{}}\PYG{l+s+s1}{utf\PYGZhy{}8}\PYG{l+s+s1}{\PYGZsq{}}\PYG{p}{)} \PYG{k}{as} \PYG{n}{plik}\PYG{p}{:}
        \PYG{n}{csv\PYGZus{}reader} \PYG{o}{=} \PYG{n}{csv}\PYG{o}{.}\PYG{n}{reader}\PYG{p}{(}\PYG{n}{plik}\PYG{p}{)}
        \PYG{n+nb}{next}\PYG{p}{(}\PYG{n}{csv\PYGZus{}reader}\PYG{p}{)}  \PYG{c+c1}{\PYGZsh{} Pominięcie wiersza nagłówka}
        \PYG{k}{for} \PYG{n}{wiersz} \PYG{o+ow}{in} \PYG{n}{csv\PYGZus{}reader}\PYG{p}{:}
            \PYG{n}{dane}\PYG{o}{.}\PYG{n}{append}\PYG{p}{(}\PYG{n+nb}{tuple}\PYG{p}{(}\PYG{n}{wiersz}\PYG{p}{)}\PYG{p}{)}
    \PYG{k}{return} \PYG{n}{dane}

\PYG{c+c1}{\PYGZsh{} 2. Właściwy proces wsadowego importu do bazy PostgreSQL}
\PYG{k}{def} \PYG{n+nf}{importuj\PYGZus{}do\PYGZus{}postgres}\PYG{p}{(}\PYG{n}{kategorie\PYGZus{}do\PYGZus{}importu}\PYG{p}{,} \PYG{n}{db\PYGZus{}config}\PYG{p}{)}\PYG{p}{:}
    \PYG{k}{try}\PYG{p}{:}
        \PYG{c+c1}{\PYGZsh{} Użycie Context Managera dla bezpiecznego zarządzania transakcjami}
        \PYG{k}{with} \PYG{n}{psycopg}\PYG{o}{.}\PYG{n}{connect}\PYG{p}{(}\PYG{o}{*}\PYG{o}{*}\PYG{n}{db\PYGZus{}config}\PYG{p}{)} \PYG{k}{as} \PYG{n}{conn}\PYG{p}{:}
            \PYG{k}{with} \PYG{n}{conn}\PYG{o}{.}\PYG{n}{cursor}\PYG{p}{(}\PYG{p}{)} \PYG{k}{as} \PYG{n}{cursor}\PYG{p}{:}

                \PYG{c+c1}{\PYGZsh{} Definicja sparametryzowanego zapytania zapobiegającego SQL Injection}
                \PYG{n}{zapytanie\PYGZus{}sql} \PYG{o}{=} \PYG{l+s+s2}{\PYGZdq{}}\PYG{l+s+s2}{INSERT INTO Kategorie (Nazwa\PYGZus{}Kategorii) VALUES (}\PYG{l+s+si}{\PYGZpc{}s}\PYG{l+s+s2}{)}\PYG{l+s+s2}{\PYGZdq{}}

                \PYG{c+c1}{\PYGZsh{} Wsadowe wywołanie zapytań (batch insert)}
                \PYG{n}{cursor}\PYG{o}{.}\PYG{n}{executemany}\PYG{p}{(}\PYG{n}{zapytanie\PYGZus{}sql}\PYG{p}{,} \PYG{n}{kategorie\PYGZus{}do\PYGZus{}importu}\PYG{p}{)}

                \PYG{n+nb}{print}\PYG{p}{(}\PYG{l+s+sa}{f}\PYG{l+s+s2}{\PYGZdq{}}\PYG{l+s+s2}{Pomyślnie zaimportowano }\PYG{l+s+si}{\PYGZob{}}\PYG{n+nb}{len}\PYG{p}{(}\PYG{n}{kategorie\PYGZus{}do\PYGZus{}importu}\PYG{p}{)}\PYG{l+s+si}{\PYGZcb{}}\PYG{l+s+s2}{ rekordów.}\PYG{l+s+s2}{\PYGZdq{}}\PYG{p}{)}

    \PYG{k}{except} \PYG{n+ne}{Exception} \PYG{k}{as} \PYG{n}{e}\PYG{p}{:}
        \PYG{n+nb}{print}\PYG{p}{(}\PYG{l+s+sa}{f}\PYG{l+s+s2}{\PYGZdq{}}\PYG{l+s+s2}{Wystąpił błąd operacji wsadowej: }\PYG{l+s+si}{\PYGZob{}}\PYG{n}{e}\PYG{l+s+si}{\PYGZcb{}}\PYG{l+s+s2}{\PYGZdq{}}\PYG{p}{)}
\end{sphinxVerbatim}


\section{Podsumowanie}
\label{\detokenize{rozdzial_4/index:podsumowanie}}
\sphinxAtStartPar
Niniejszy rozdział wieńczy etap projektowania i wdrażania struktury relacyjnej bazy danych dla systemu zarządzania biblioteką. Poprzez transformację modelu logicznego do postaci fizycznej osiągnięto kompletne i funkcjonalne rozwiązanie.

\sphinxAtStartPar
Pierwsza część rozdziału koncentrowała się na \sphinxstylestrong{Definicji fizycznej bazy danych (Zadanie 1)}, gdzie dokonano implementacji schematów dla dwóch wiodących silników relacyjnych: PostgreSQL i SQLite, stanowiących praktyczną alternatywę w zależności od wymagań skalowalności oraz zasobów dostępnych w środowisku produkcyjnym.

\sphinxAtStartPar
Druga część rozdziału poświęcona została \sphinxstylestrong{Mechanizmowi zasilania bazy danych (Zadanie 2)}. Dane demonstracyjne przygotowano w formacie CSV, który zapewnia elastyczność i łatwość edycji zarówno dla administratorów baz danych, jak i dla użytkowników końcowych systemu.

\sphinxAtStartPar
Automatyzacja procesu inicjalizacji systemu za pomocą skryptów Python stanowi znaczący wkład w praktyczność rozwiązania, umożliwiając łatwe replikowanie i testowanie bazy danych w różnych środowiskach bez konieczności wykonywania repetycyjnych czynności manualnych.

\sphinxAtStartPar
Niniejszy projekt stanowi kompletne studium cyklu życia relacyjnej bazy danych \textendash{} od teoretycznego modelowania (Rozdziały 1\sphinxhyphen{}3) poprzez praktyczną implementację (Rozdział 4) \textendash{} demonstrując całkowity proces od koncepcji do produkcyjnego systemu zarządzania danymi.

\sphinxstepscope


\chapter{Rozdział 5: Komunikacja i operacje na danych}
\label{\detokenize{rozdzial_5/index:rozdzial-5-komunikacja-i-operacje-na-danych}}\label{\detokenize{rozdzial_5/index::doc}}\begin{quote}\begin{description}
\sphinxlineitem{Autorzy}\begin{enumerate}
\sphinxsetlistlabels{\arabic}{enumi}{enumii}{}{.}%
\item {} 
\sphinxAtStartPar
Paweł Łoćwin

\item {} 
\sphinxAtStartPar
Paweł Łosowski

\end{enumerate}

\end{description}\end{quote}


\section{Wstęp}
\label{\detokenize{rozdzial_5/index:wstep}}
\sphinxAtStartPar
Po zdefiniowaniu i zasileniu struktur bazodanowych w poprzednich rozdziałach, system wymaga
interfejsu zdolnego do obsługi zaawansowanej logiki biznesowej biblioteki. W tym celu zaimplementowano
warstwę dostępu do danych (Data Access Layer) w języku Python.

\sphinxAtStartPar
Wykorzystano techniki tworzenia rozbudowanych zapytań SQL, w tym:
* Podzapytania (skorelowane i nieskorelowane w klauzulach \sphinxcode{\sphinxupquote{SELECT}}, \sphinxcode{\sphinxupquote{FROM}}, \sphinxcode{\sphinxupquote{WHERE}}).
* Złączenia relacyjne (\sphinxcode{\sphinxupquote{JOIN}}, \sphinxcode{\sphinxupquote{LEFT JOIN}}).
* Analizę agregacyjną (\sphinxcode{\sphinxupquote{GROUP BY}}, \sphinxcode{\sphinxupquote{HAVING}}).
* Operatory teorii mnogości (\sphinxcode{\sphinxupquote{UNION}}, \sphinxcode{\sphinxupquote{INTERSECT}}, \sphinxcode{\sphinxupquote{EXCEPT}}).


\section{Pełna specyfikacja stworzonego API}
\label{\detokenize{rozdzial_5/index:pelna-specyfikacja-stworzonego-api}}
\sphinxAtStartPar
Poniżej przedstawiono zautomatyzowaną dokumentację modułu obsługującego połączenia z
bazami PostgreSQL oraz SQLite, wygenerowaną na podstawie wbudowanych komentarzy (Docstrings).
\index{module@\spxentry{module}!biblioteka\_zapytania@\spxentry{biblioteka\_zapytania}}\index{biblioteka\_zapytania@\spxentry{biblioteka\_zapytania}!module@\spxentry{module}}

\subsection{Moduł biblioteka\_zapytania}
\label{\detokenize{rozdzial_5/index:modul-biblioteka-zapytania}}\label{\detokenize{rozdzial_5/index:module-biblioteka_zapytania}}
\sphinxAtStartPar
Moduł zawiera zestaw zaawansowanych zapytań SQL do obsługi systemu zarządzania biblioteką.
Obsługuje dwa silniki bazodanowe: PostgreSQL oraz SQLite.

\sphinxAtStartPar
Wymagania:
\sphinxhyphen{} psycopg2 (lub psycopg) dla PostgreSQL
\sphinxhyphen{} sqlite3 (wbudowany) dla SQLite
\index{get\_aktywni\_czytelnicy\_powyzej\_limitu() (in module biblioteka\_zapytania)@\spxentry{get\_aktywni\_czytelnicy\_powyzej\_limitu()}\spxextra{in module biblioteka\_zapytania}}

\begin{savenotes}\begin{fulllineitems}
\phantomsection\label{\detokenize{rozdzial_5/index:biblioteka_zapytania.get_aktywni_czytelnicy_powyzej_limitu}}
\pysigstartsignatures
\pysiglinewithargsret
{\sphinxcode{\sphinxupquote{biblioteka\_zapytania.}}\sphinxbfcode{\sphinxupquote{get\_aktywni\_czytelnicy\_powyzej\_limitu}}}
{\sphinxparam{\DUrole{n}{conn}}\sphinxparamcomma \sphinxparam{\DUrole{n}{limit\_wypozyczen}}}
{}
\pysigstopsignatures
\sphinxAtStartPar
Pobiera listę czytelników, którzy w historii wypożyczyli więcej książek niż wynosi podany limit.
Wykorzystuje funkcje agregujące, złączenia oraz klauzulę HAVING.
\begin{quote}\begin{description}
\sphinxlineitem{Parameters}\begin{itemize}
\item {} 
\sphinxAtStartPar
\sphinxstyleliteralstrong{\sphinxupquote{conn}} (\sphinxstyleliteralemphasis{\sphinxupquote{psycopg2.extensions.connection}}) \textendash{} Obiekt aktywnego połączenia z bazą PostgreSQL (psycopg2).

\item {} 
\sphinxAtStartPar
\sphinxstyleliteralstrong{\sphinxupquote{limit\_wypozyczen}} (\sphinxstyleliteralemphasis{\sphinxupquote{int}}) \textendash{} Minimalna liczba wypożyczeń uprawniająca do znalezienia się na liście.

\end{itemize}

\sphinxlineitem{Returns}
\sphinxAtStartPar
Lista krotek zawierająca imię, nazwisko i sumę wypożyczeń czytelnika.

\sphinxlineitem{Return type}
\sphinxAtStartPar
list{[}tuple{]}

\end{description}\end{quote}

\end{fulllineitems}\end{savenotes}

\index{get\_czytelnicy\_bez\_wypozyczen() (in module biblioteka\_zapytania)@\spxentry{get\_czytelnicy\_bez\_wypozyczen()}\spxextra{in module biblioteka\_zapytania}}

\begin{savenotes}\begin{fulllineitems}
\phantomsection\label{\detokenize{rozdzial_5/index:biblioteka_zapytania.get_czytelnicy_bez_wypozyczen}}
\pysigstartsignatures
\pysiglinewithargsret
{\sphinxcode{\sphinxupquote{biblioteka\_zapytania.}}\sphinxbfcode{\sphinxupquote{get\_czytelnicy\_bez\_wypozyczen}}}
{\sphinxparam{\DUrole{n}{conn}}}
{}
\pysigstopsignatures
\sphinxAtStartPar
Pobiera dane czytelników, którzy zapisali się do biblioteki, ale nigdy nie wypożyczyli żadnej książki.
Zastosowano tu operator zbiorowy EXCEPT w celu odjęcia zbioru aktywnych od wszystkich.
\begin{quote}\begin{description}
\sphinxlineitem{Parameters}
\sphinxAtStartPar
\sphinxstyleliteralstrong{\sphinxupquote{conn}} (\sphinxstyleliteralemphasis{\sphinxupquote{psycopg2.extensions.connection}}) \textendash{} Obiekt aktywnego połączenia z bazą PostgreSQL.

\sphinxlineitem{Returns}
\sphinxAtStartPar
Lista krotek z danymi czytelników (ID, Imię, Nazwisko).

\sphinxlineitem{Return type}
\sphinxAtStartPar
list{[}tuple{]}

\end{description}\end{quote}

\end{fulllineitems}\end{savenotes}

\index{get\_katalog\_osob() (in module biblioteka\_zapytania)@\spxentry{get\_katalog\_osob()}\spxextra{in module biblioteka\_zapytania}}

\begin{savenotes}\begin{fulllineitems}
\phantomsection\label{\detokenize{rozdzial_5/index:biblioteka_zapytania.get_katalog_osob}}
\pysigstartsignatures
\pysiglinewithargsret
{\sphinxcode{\sphinxupquote{biblioteka\_zapytania.}}\sphinxbfcode{\sphinxupquote{get\_katalog\_osob}}}
{\sphinxparam{\DUrole{n}{conn}}}
{}
\pysigstopsignatures
\sphinxAtStartPar
Zwraca ujednoliconą listę wszystkich osób figurujących w systemie, przypisując im odpowiednią
rolę. Używa operatora zbiorowego UNION w celu złączenia tabel Czytelnicy i Autorzy.
\begin{quote}\begin{description}
\sphinxlineitem{Parameters}
\sphinxAtStartPar
\sphinxstyleliteralstrong{\sphinxupquote{conn}} (\sphinxstyleliteralemphasis{\sphinxupquote{sqlite3.Connection}}) \textendash{} Obiekt połączenia z bazą SQLite.

\sphinxlineitem{Returns}
\sphinxAtStartPar
Lista krotek postaci (Imię, Nazwisko, Rola).

\sphinxlineitem{Return type}
\sphinxAtStartPar
list{[}tuple{]}

\end{description}\end{quote}

\end{fulllineitems}\end{savenotes}

\index{get\_ksiazki\_wypozyczone\_w\_miescie() (in module biblioteka\_zapytania)@\spxentry{get\_ksiazki\_wypozyczone\_w\_miescie()}\spxextra{in module biblioteka\_zapytania}}

\begin{savenotes}\begin{fulllineitems}
\phantomsection\label{\detokenize{rozdzial_5/index:biblioteka_zapytania.get_ksiazki_wypozyczone_w_miescie}}
\pysigstartsignatures
\pysiglinewithargsret
{\sphinxcode{\sphinxupquote{biblioteka\_zapytania.}}\sphinxbfcode{\sphinxupquote{get\_ksiazki\_wypozyczone\_w\_miescie}}}
{\sphinxparam{\DUrole{n}{conn}}\sphinxparamcomma \sphinxparam{\DUrole{n}{nazwa\_miasta}}}
{}
\pysigstopsignatures
\sphinxAtStartPar
Wyszukuje tytuły książek, które kiedykolwiek zostały wypożyczone przez mieszkańców określonego miasta.
Rozwiązanie bazuje na podzapytaniu skorelowanym w klauzuli WHERE.
\begin{quote}\begin{description}
\sphinxlineitem{Parameters}\begin{itemize}
\item {} 
\sphinxAtStartPar
\sphinxstyleliteralstrong{\sphinxupquote{conn}} (\sphinxstyleliteralemphasis{\sphinxupquote{psycopg2.extensions.connection}}) \textendash{} Obiekt aktywnego połączenia z bazą PostgreSQL.

\item {} 
\sphinxAtStartPar
\sphinxstyleliteralstrong{\sphinxupquote{nazwa\_miasta}} (\sphinxstyleliteralemphasis{\sphinxupquote{str}}) \textendash{} Nazwa miasta do przefiltrowania czytelników.

\end{itemize}

\sphinxlineitem{Returns}
\sphinxAtStartPar
Lista krotek zawierająca wyłącznie tytuły książek.

\sphinxlineitem{Return type}
\sphinxAtStartPar
list{[}tuple{]}

\end{description}\end{quote}

\end{fulllineitems}\end{savenotes}

\index{get\_ksiazki\_z\_iloscia\_wypozyczen() (in module biblioteka\_zapytania)@\spxentry{get\_ksiazki\_z\_iloscia\_wypozyczen()}\spxextra{in module biblioteka\_zapytania}}

\begin{savenotes}\begin{fulllineitems}
\phantomsection\label{\detokenize{rozdzial_5/index:biblioteka_zapytania.get_ksiazki_z_iloscia_wypozyczen}}
\pysigstartsignatures
\pysiglinewithargsret
{\sphinxcode{\sphinxupquote{biblioteka\_zapytania.}}\sphinxbfcode{\sphinxupquote{get\_ksiazki\_z\_iloscia\_wypozyczen}}}
{\sphinxparam{\DUrole{n}{conn}}\sphinxparamcomma \sphinxparam{\DUrole{n}{rok\_od}}}
{}
\pysigstopsignatures
\sphinxAtStartPar
Zwraca szczegółowe dane o książkach wydanych po zadanym roku wraz z ilością ich historycznych 
wypożyczeń. Wykorzystuje funkcje wierszowe (UPPER) oraz podzapytanie w klauzuli SELECT.
\begin{quote}\begin{description}
\sphinxlineitem{Parameters}\begin{itemize}
\item {} 
\sphinxAtStartPar
\sphinxstyleliteralstrong{\sphinxupquote{conn}} (\sphinxstyleliteralemphasis{\sphinxupquote{psycopg2.extensions.connection}}) \textendash{} Obiekt aktywnego połączenia z bazą PostgreSQL.

\item {} 
\sphinxAtStartPar
\sphinxstyleliteralstrong{\sphinxupquote{rok\_od}} (\sphinxstyleliteralemphasis{\sphinxupquote{int}}) \textendash{} Dolna granica roku wydania książki (wyłącznie).

\end{itemize}

\sphinxlineitem{Returns}
\sphinxAtStartPar
Lista krotek (Tytuł, ZWielkiejLitery\_NazwiskoAutora, Total\_Wypozyczen).

\sphinxlineitem{Return type}
\sphinxAtStartPar
list{[}tuple{]}

\end{description}\end{quote}

\end{fulllineitems}\end{savenotes}

\index{get\_najdluzej\_przetrzymywane() (in module biblioteka\_zapytania)@\spxentry{get\_najdluzej\_przetrzymywane()}\spxextra{in module biblioteka\_zapytania}}

\begin{savenotes}\begin{fulllineitems}
\phantomsection\label{\detokenize{rozdzial_5/index:biblioteka_zapytania.get_najdluzej_przetrzymywane}}
\pysigstartsignatures
\pysiglinewithargsret
{\sphinxcode{\sphinxupquote{biblioteka\_zapytania.}}\sphinxbfcode{\sphinxupquote{get\_najdluzej\_przetrzymywane}}}
{\sphinxparam{\DUrole{n}{conn}}\sphinxparamcomma \sphinxparam{\DUrole{n}{limit\_rekordow}}}
{}
\pysigstopsignatures
\sphinxAtStartPar
Wyszukuje wypożyczenia, które aktualnie trwają (brak daty zwrotu) i sortuje je rosnąco od 
najstarszych. Wykorzystuje funkcję wierszową DATE() silnika SQLite działającą na ciągach TEXT.
\begin{quote}\begin{description}
\sphinxlineitem{Parameters}\begin{itemize}
\item {} 
\sphinxAtStartPar
\sphinxstyleliteralstrong{\sphinxupquote{conn}} (\sphinxstyleliteralemphasis{\sphinxupquote{sqlite3.Connection}}) \textendash{} Obiekt połączenia z bazą SQLite.

\item {} 
\sphinxAtStartPar
\sphinxstyleliteralstrong{\sphinxupquote{limit\_rekordow}} (\sphinxstyleliteralemphasis{\sphinxupquote{int}}) \textendash{} Liczba rekordów do zwrócenia (LIMIT).

\end{itemize}

\sphinxlineitem{Returns}
\sphinxAtStartPar
Lista krotek (Tytuł, Imię, Nazwisko, Data wypożyczenia).

\sphinxlineitem{Return type}
\sphinxAtStartPar
list{[}tuple{]}

\end{description}\end{quote}

\end{fulllineitems}\end{savenotes}

\index{get\_puste\_kategorie() (in module biblioteka\_zapytania)@\spxentry{get\_puste\_kategorie()}\spxextra{in module biblioteka\_zapytania}}

\begin{savenotes}\begin{fulllineitems}
\phantomsection\label{\detokenize{rozdzial_5/index:biblioteka_zapytania.get_puste_kategorie}}
\pysigstartsignatures
\pysiglinewithargsret
{\sphinxcode{\sphinxupquote{biblioteka\_zapytania.}}\sphinxbfcode{\sphinxupquote{get\_puste\_kategorie}}}
{\sphinxparam{\DUrole{n}{conn}}}
{}
\pysigstopsignatures
\sphinxAtStartPar
Pobiera nazwy kategorii, do których aktualnie nie przypisano żadnych książek w inwentarzu.
Demonstruje wykorzystanie LEFT JOIN w celu znalezienia brakujących relacji (NULL).
\begin{quote}\begin{description}
\sphinxlineitem{Parameters}
\sphinxAtStartPar
\sphinxstyleliteralstrong{\sphinxupquote{conn}} (\sphinxstyleliteralemphasis{\sphinxupquote{sqlite3.Connection}}) \textendash{} Obiekt połączenia z bazą SQLite.

\sphinxlineitem{Returns}
\sphinxAtStartPar
Lista krotek z nazwami “pustych” kategorii.

\sphinxlineitem{Return type}
\sphinxAtStartPar
list{[}tuple{]}

\end{description}\end{quote}

\end{fulllineitems}\end{savenotes}

\index{get\_raport\_autorow() (in module biblioteka\_zapytania)@\spxentry{get\_raport\_autorow()}\spxextra{in module biblioteka\_zapytania}}

\begin{savenotes}\begin{fulllineitems}
\phantomsection\label{\detokenize{rozdzial_5/index:biblioteka_zapytania.get_raport_autorow}}
\pysigstartsignatures
\pysiglinewithargsret
{\sphinxcode{\sphinxupquote{biblioteka\_zapytania.}}\sphinxbfcode{\sphinxupquote{get\_raport\_autorow}}}
{\sphinxparam{\DUrole{n}{conn}}}
{}
\pysigstopsignatures
\sphinxAtStartPar
Tworzy zaawansowany raport o autorach, zliczając ich wszystkie książki oraz wskazując 
tytuł chronologicznie najnowszego dzieła za pomocą podzapytania w SELECT.
\begin{quote}\begin{description}
\sphinxlineitem{Parameters}
\sphinxAtStartPar
\sphinxstyleliteralstrong{\sphinxupquote{conn}} (\sphinxstyleliteralemphasis{\sphinxupquote{sqlite3.Connection}}) \textendash{} Obiekt połączenia z bazą SQLite.

\sphinxlineitem{Returns}
\sphinxAtStartPar
Lista krotek (Imię autora, Nazwisko, Suma dzieł, Tytuł najnowszej książki).

\sphinxlineitem{Return type}
\sphinxAtStartPar
list{[}tuple{]}

\end{description}\end{quote}

\end{fulllineitems}\end{savenotes}

\index{get\_srednia\_wypozyczen\_na\_czytelnika() (in module biblioteka\_zapytania)@\spxentry{get\_srednia\_wypozyczen\_na\_czytelnika()}\spxextra{in module biblioteka\_zapytania}}

\begin{savenotes}\begin{fulllineitems}
\phantomsection\label{\detokenize{rozdzial_5/index:biblioteka_zapytania.get_srednia_wypozyczen_na_czytelnika}}
\pysigstartsignatures
\pysiglinewithargsret
{\sphinxcode{\sphinxupquote{biblioteka\_zapytania.}}\sphinxbfcode{\sphinxupquote{get\_srednia\_wypozyczen\_na\_czytelnika}}}
{\sphinxparam{\DUrole{n}{conn}}}
{}
\pysigstopsignatures
\sphinxAtStartPar
Oblicza średnią liczbę wypożyczonych książek przypadającą na jednego aktywnego czytelnika.
Wykorzystuje podzapytanie w klauzuli FROM do wcześniejszej agregacji danych.
\begin{quote}\begin{description}
\sphinxlineitem{Parameters}
\sphinxAtStartPar
\sphinxstyleliteralstrong{\sphinxupquote{conn}} (\sphinxstyleliteralemphasis{\sphinxupquote{sqlite3.Connection}}) \textendash{} Obiekt połączenia z bazą SQLite.

\sphinxlineitem{Returns}
\sphinxAtStartPar
Zwraca jedną krotkę z pojedynczą wartością zmiennoprzecinkową (średnia).

\sphinxlineitem{Return type}
\sphinxAtStartPar
list{[}tuple{]}

\end{description}\end{quote}

\end{fulllineitems}\end{savenotes}

\index{get\_wszechstronni\_czytelnicy() (in module biblioteka\_zapytania)@\spxentry{get\_wszechstronni\_czytelnicy()}\spxextra{in module biblioteka\_zapytania}}

\begin{savenotes}\begin{fulllineitems}
\phantomsection\label{\detokenize{rozdzial_5/index:biblioteka_zapytania.get_wszechstronni_czytelnicy}}
\pysigstartsignatures
\pysiglinewithargsret
{\sphinxcode{\sphinxupquote{biblioteka\_zapytania.}}\sphinxbfcode{\sphinxupquote{get\_wszechstronni\_czytelnicy}}}
{\sphinxparam{\DUrole{n}{conn}}}
{}
\pysigstopsignatures
\sphinxAtStartPar
Zwraca czytelników (Imię i Nazwisko), którzy wypożyczali książki z kategorii ‘Fantastyka’
oraz jednocześnie z kategorii ‘Kryminał’. Używa operatora zbiorowego INTERSECT.
\begin{quote}\begin{description}
\sphinxlineitem{Parameters}
\sphinxAtStartPar
\sphinxstyleliteralstrong{\sphinxupquote{conn}} (\sphinxstyleliteralemphasis{\sphinxupquote{psycopg2.extensions.connection}}) \textendash{} Obiekt aktywnego połączenia z bazą PostgreSQL.

\sphinxlineitem{Returns}
\sphinxAtStartPar
Lista krotek z imionami i nazwiskami wszechstronnych czytelników.

\sphinxlineitem{Return type}
\sphinxAtStartPar
list{[}tuple{]}

\end{description}\end{quote}

\end{fulllineitems}\end{savenotes}




\renewcommand{\indexname}{Index}

\end{document}